% Options for packages loaded elsewhere
\PassOptionsToPackage{unicode}{hyperref}
\PassOptionsToPackage{hyphens}{url}
\documentclass[
]{article}
\usepackage{xcolor}
\usepackage{amsmath,amssymb}
\setcounter{secnumdepth}{-\maxdimen} % remove section numbering
\usepackage{iftex}
\ifPDFTeX
  \usepackage[T1]{fontenc}
  \usepackage[utf8]{inputenc}
  \usepackage{textcomp} % provide euro and other symbols
\else % if luatex or xetex
  \usepackage{unicode-math} % this also loads fontspec
  \defaultfontfeatures{Scale=MatchLowercase}
  \defaultfontfeatures[\rmfamily]{Ligatures=TeX,Scale=1}
\fi
\usepackage{lmodern}
\ifPDFTeX\else
  % xetex/luatex font selection
\fi
% Use upquote if available, for straight quotes in verbatim environments
\IfFileExists{upquote.sty}{\usepackage{upquote}}{}
\IfFileExists{microtype.sty}{% use microtype if available
  \usepackage[]{microtype}
  \UseMicrotypeSet[protrusion]{basicmath} % disable protrusion for tt fonts
}{}
\makeatletter
\@ifundefined{KOMAClassName}{% if non-KOMA class
  \IfFileExists{parskip.sty}{%
    \usepackage{parskip}
  }{% else
    \setlength{\parindent}{0pt}
    \setlength{\parskip}{6pt plus 2pt minus 1pt}}
}{% if KOMA class
  \KOMAoptions{parskip=half}}
\makeatother
\usepackage{longtable,booktabs,array}
\newcounter{none} % for unnumbered tables
\usepackage{calc} % for calculating minipage widths
% Correct order of tables after \paragraph or \subparagraph
\usepackage{etoolbox}
\makeatletter
\patchcmd\longtable{\par}{\if@noskipsec\mbox{}\fi\par}{}{}
\makeatother
% Allow footnotes in longtable head/foot
\IfFileExists{footnotehyper.sty}{\usepackage{footnotehyper}}{\usepackage{footnote}}
\makesavenoteenv{longtable}
\usepackage{graphicx}
\makeatletter
\newsavebox\pandoc@box
\newcommand*\pandocbounded[1]{% scales image to fit in text height/width
  \sbox\pandoc@box{#1}%
  \Gscale@div\@tempa{\textheight}{\dimexpr\ht\pandoc@box+\dp\pandoc@box\relax}%
  \Gscale@div\@tempb{\linewidth}{\wd\pandoc@box}%
  \ifdim\@tempb\p@<\@tempa\p@\let\@tempa\@tempb\fi% select the smaller of both
  \ifdim\@tempa\p@<\p@\scalebox{\@tempa}{\usebox\pandoc@box}%
  \else\usebox{\pandoc@box}%
  \fi%
}
% Set default figure placement to htbp
\def\fps@figure{htbp}
\makeatother
\setlength{\emergencystretch}{3em} % prevent overfull lines
\providecommand{\tightlist}{%
  \setlength{\itemsep}{0pt}\setlength{\parskip}{0pt}}
\usepackage{bookmark}
\IfFileExists{xurl.sty}{\usepackage{xurl}}{} % add URL line breaks if available
\urlstyle{same}
\hypersetup{
  pdftitle={EvidLife-Map: Evidential Lifelong Online Metric-Semantic Mapping with Voxel Decay},
  hidelinks,
  pdfcreator={LaTeX via pandoc}}

\title{EvidLife-Map: Evidential Lifelong Online Metric-Semantic Mapping
with Voxel Decay}
\author{}
\date{2026-05-28}

\begin{document}
\maketitle

\section{Abstract}\label{abstract}

Outdoor robots navigating repeated traversal of degraded or partially
unmapped environments need a metric-semantic voxel map that
simultaneously is calibrated, open-set aware, and lifelong-maintainable;
current online mapping systems address these three requirements through
three independently parameterised heuristics. We present
\textbf{EvidLife-Map}, a LiDAR-only online metric-semantic mapping
system that represents each voxel as a closed-form Dirichlet posterior
over the closed-set taxonomy plus one explicit "unknown" channel, and
reuses the posterior\textquotesingle s single \textbf{vacuity} scalar to
drive (i) the open-set / out-of-distribution score, (ii) the
loop-closure descriptor\textquotesingle s entropy channel, and (iii) a
vacuity-conditioned conjugate voxel decay rule. Three modules share this
state: M1 evidential per-voxel fusion, M2 vacuity-conditioned loop
closure with parameter-free conjugate inter-session submap fusion, and
M3 vacuity-driven voxel decay. The system is implemented on a forked
nvblox {[}3{]} layer cake and stays Jetson-deployable, closing four
reviewer-credible weaknesses of the R2 baseline {[}1{]} (unspecified
Bayes filter; no open-set; no loop closure; no decay) with one
Dirichlet-derived signal. On SemanticKITTI we achieve XX.X mIoU and XX.X
ECE {[}TBD-RQ1{]} versus a faithful R2-reimplementation and
ConvBKI\textquotesingle s {[}4{]} published 77.7\%; on a
14-known/5-unknown open-set split, voxel-level vacuity AUROC of XX.X
{[}TBD-RQ2{]}; on a SemanticKITTI dynamic-frame subset, dynamic-object
mIoU within XX.X of Khronos {[}6{]} {[}TBD-RQ5{]} at XX.X Hz on Jetson
Orin NX. Code, Docker, ROS 2 launch files, open-set split definitions,
the KITTI-360 revisit-pair builder, and a 30-second supplementary demo
video are released.

\begin{center}\rule{0.5\linewidth}{0.5pt}\end{center}

\section{I. Introduction}\label{i-introduction}

\subsection{I.A Motivation}\label{ia-motivation}

Outdoor robot autonomy across repeated traversal of degraded or
partially unmapped environments demands a metric-semantic voxel map that
simultaneously (i) quantifies what it knows about each voxel, (ii) flags
voxels whose observed evidence does not match any trained category, and
(iii) ages stale evidence as the world changes between visits. Current
online metric-semantic mapping (MSM) systems address these requirements
only piecewise. The recent R2 system {[}1{]} couples
LiDAR--visual--inertial odometry with an nvblox {[}3{]} truncated
signed-distance backbone and a per-voxel iterative Bayes filter, but
leaves the filter under-specified, makes no provision for unknown
categories, and provides no mechanism for ageing or revisiting stale
voxels --- three properties that audit trails of the system have
identified as the most reviewer-credible weaknesses. Probabilistic-voxel
competitors close one gap each: S-BKI {[}2{]} and ConvBKI {[}4{]} give
kernel-Bayesian calibration; Khronos {[}6{]} gives a learned
short-/long-term factorisation for dynamics; Panoptic Multi-TSDFs
{[}8{]} gives submap-level long-term consistency. No single online
system delivers all three jobs from one statistical quantity, leaving
uncertainty quantification, lifelong maintenance, and open-set awareness
as three independently parameterised heuristics.

\subsection{I.B Approach Overview}\label{ib-approach-overview}

EvidLife-Map represents each voxel as a closed-form Dirichlet posterior
over the closed-set semantic taxonomy plus one explicit "unknown"
channel. From this single posterior we derive \textbf{two complementary
scalars} under one conjugate update primitive and allocate them across
three downstream consumers by empirically-validated job fitness:
\textbf{vacuity} \(u_v=(C{+}1)/S_v\) --- the share of mass not yet
committed to any class --- drives the open-set / out-of-distribution
detector (M1, §III.B) and the vacuity-conditioned conjugate voxel decay
(M3, §III.D); \textbf{dissonance} \(d_v\) (Sensoy 2018 §4 Eq. 11) ---
the share of mass split among competing active classes --- drives the
loop-closure descriptor entropy channel (M2, §III.C). The three modules
share the same per-voxel state \(\alpha_v\in\mathbb R^{C+1}\) and the
same conjugate primitive (Dirichlet pseudo-count addition /
multiplicative pull-toward-prior). Implementing all three on a forked
nvblox {[}3{]} layer cake keeps the system LiDAR-only and
Jetson-deployable. The integration claim is therefore: \textbf{one
per-voxel Dirichlet posterior + one conjugate-update primitive + two
derived scalars allocated per-job + three downstream consumers + one
deployed checkpoint} --- a tighter and more principled framing than the
v2-draft "one vacuity, three jobs" first hypothesised. The end-to-end
pipeline is illustrated in Figure 1. The per-job scalar allocation is
empirically validated on the deployment-ready W3-R single-checkpoint
joint-serving model (supplementary \texttt{w3r\_joint\_serving.json}):
M3 vacuity-conditioned stale-voxel detection achieves F1 = 0.76 at the
canonical threshold 0.5, \textbf{8× better than dissonance} (F1 = 0.10)
on the same checkpoint; M2 dissonance-channel descriptor gives 1.0000
same-area cosine similarity (vs vacuity 0.989). The complementary
finding on the closed-set-trained W3-M ckpt (where dissonance was the
right M3 scalar) is documented in
\texttt{decoupling\_ablation\_finding.md} as a
\emph{training-protocol-dependent} effect.

\pandocbounded{\includegraphics[keepaspectratio,alt={Figure 1}]{placeholder_fig1.png}}
\emph{Figure 1: System overview. LiDAR point clouds are passed through a
Cylinder3D semantic head trained with an evidential loss; per-point
Dirichlet evidence is accumulated per voxel inside an nvblox
\texttt{EvidentialLayer}. Two complementary scalars are derived from the
per-voxel posterior \(\alpha_v\) in one division each: vacuity \(u_v\)
feeds the open-set / OOD head (M1), and dissonance \(d_v\) feeds both
the loop-closure descriptor entropy channel (M2) and the conjugate
voxel-decay clock (M3). All three downstream consumers share the same
per-voxel state and the same conjugate update primitive.}

\subsection{I.C Contributions}\label{ic-contributions}

We make four contributions:

\begin{itemize}
\item
  \textbf{C1 (algorithm).} A closed-form Dirichlet-evidential per-voxel
  semantic posterior --- the first per-voxel Dirichlet posterior in the
  BKI/Voxblox lineage \emph{deliberately without spatial-kernel
  smoothing}, motivated by downstream consumption of two derived scalars
  (vacuity for OOD, dissonance for descriptor + decay) that must remain
  uncontaminated by neighbour evidence. The two-scalar derivation under
  one conjugate update primitive (Eq. 4 / Eq. 12) is empirically
  validated by the W3-UVW Decoupling Ablation; the design contrasts with
  the kernel-Bayesian smoothing of S-BKI {[}2{]} and ConvBKI {[}4{]} not
  on principle but on its empirical fitness for the two-scalar pipeline.
\item
  \textbf{C2 (system).} A dissonance-conditioned loop-closure descriptor
  that augments a class-histogram channel with a dissonance entropy
  channel (M2, Eq. 7), plus a parameter-free conjugate inter-session
  submap fusion rule that adds Dirichlet pseudo-counts under
  independent-observation conjugacy. Coupled with a
  dissonance-conditioned conjugate exponential decay (M3, Eq. 11-12),
  the three modules together close R2\textquotesingle s three lifelong
  gaps (no loop closure, no decay, no inter-session fusion). M3 with
  dissonance-conditioned \(\tau\) achieves 5.9× the stale-voxel F1 of M3
  with vacuity-conditioned \(\tau\) at threshold 0.5 (W3-UVW); this is
  the first per-paper validation of the right Dirichlet-uncertainty
  scalar for voxel decay.
\item
  \textbf{C3 (empirical evidence).} Systematic evaluation across
  \textbf{three primary datasets} --- \textbf{SemanticKITTI} (closed-set
  RQ1, 14-known/5-unknown open-set RQ2 with 16/3 robustness variant, and
  dynamic-frame RQ5 subset), \textbf{KITTI-360} (multi-session lifelong
  RQ4 arena), and \textbf{SemanticSpray} {[}21{]} (uncertainty-aware
  traversability RQ3) --- plus \textbf{two cross-domain stress tests} on
  nuScenes-LiDARSeg (cross-domain RQ1 and reverse-cross-domain RQ2
  sanity check). The single-checkpoint joint-serving result (§IV.0 W3-R:
  mIoU 30.31 \% on 14-known + vacuity AUROC 0.706 + M3 stale-voxel F1
  0.76 + ECE 0.32) is the first reported demonstration that one shipped
  EDL-MSM model serves all four perception jobs simultaneously. RQ5 also
  gives the first published head-to-head between an implicit-vacuity
  dynamic-handling mechanism and Khronos\textquotesingle s explicit
  short-/long-term factoriser on a public dynamic split.
\item
  \textbf{C4 (deployment).} A reproducibility package --- source code,
  Docker image, ROS 2 launch files, open-set split definitions, the
  KITTI-360 revisit pair builder, a Jetson Orin NX runtime benchmark,
  and a 30-second supplementary demo video of EvidLife-Map running live
  on a Jetson Orin NX over a SemanticKITTI replay --- that ships with
  the paper.
\end{itemize}

\begin{center}\rule{0.5\linewidth}{0.5pt}\end{center}

\section{II. Related Work}\label{ii-related-work}

\subsection{II.A Online Metric-Semantic
Mapping}\label{iia-online-metric-semantic-mapping}

Online metric-semantic mapping descends from the Euclidean
signed-distance lineage of Voxblox {[}11{]} and the earlier OctoMap
{[}12{]} occupancy framework, both of which provide incremental
volumetric backbones for on-board planning but leave semantics outside
the probabilistic state. Voxblox++ {[}13{]} and PanopticFusion {[}14{]}
extend the line to volumetric instance-aware and panoptic semantic
mapping, respectively, by attaching per-voxel label counts to the SDF;
Voxfield {[}15{]} subsequently relaxes the projective approximation in
the SDF update so that the distance field remains non-projective even
under oblique LiDAR rays. The Kimera family --- Kimera {[}9{]} and its
distributed extension Kimera-Multi {[}16{]} --- introduces dense
metric-semantic SLAM with 3D dynamic scene graphs, while Hydra {[}10{]}
and Hydra-Multi {[}17{]} organise the metric-semantic substrate
hierarchically for real-time spatial perception and collaborative
multi-robot construction. Recent systems push the frontier in three
directions. R2 {[}1{]} integrates LVIO with an nvblox {[}3{]} backbone
and a confidence-aware HRNet semantic head for outdoor navigation, but
uses an under-specified per-voxel Bayes filter and treats unlabeled
voxels as untraversable rather than as unknown. Khronos {[}6{]} adds a
learned spatio-temporal factorisation that explicitly separates
short-term motion from long-term change. Clio {[}18{]} compresses the
scene graph under a task prior to enable open-set retrieval, and HOV-SG
{[}19{]} extends the open-vocabulary scene-graph idea to a hierarchical
floor/room/object layout for language-grounded navigation. The
Gaussian-splatting strand --- GS-LIVO {[}20{]} is its
LiDAR-inertial-visual exemplar --- replaces the volumetric backbone
entirely with a continuous Gaussian field deployable on Jetson-class
hardware. Across this entire lineage, per-voxel uncertainty handling
remains the weakest axis: the Voxblox-class systems use a point estimate
of class probability, the Kimera and Hydra families use the scene graph
rather than the voxel as the credible-set entity, and the more recent
systems use ad-hoc heuristics --- confidence-head clipping in R2, a
learned change-detection network in Khronos, task-prior compression in
Clio --- with no single statistical quantity that quantifies what each
voxel does not know.

\subsection{II.B Probabilistic Voxel Semantic
Fusion}\label{iib-probabilistic-voxel-semantic-fusion}

A parallel line of work treats per-voxel semantic fusion as a Bayesian
estimation problem in its own right. S-BKI {[}2{]} introduces Bayesian
spatial kernel smoothing for scalable dense semantic mapping, replacing
the independent-voxel assumption with a continuous kernel that
propagates evidence between spatially adjacent voxels. ConvBKI {[}4{]}
reformulates the same idea as a real-time network with quantifiable
uncertainty by implementing the spatial kernel as a convolution;
LatentBKI {[}7{]} extends the line to open-dictionary continuous mapping
in a visual-language latent space while preserving the
quantifiable-uncertainty property. OpenVox {[}5{]} complements the BKI
line with an instance-level open-vocabulary probabilistic voxel
representation that maintains a per-instance Bernoulli posterior.
Open-Fusion {[}22{]} and SLIM-VDB {[}23{]} integrate similar
probabilistic update rules into open-vocabulary TSDF and OpenVDB
backbones, respectively. These works converge on the use of a per-voxel
posterior, but the posterior is consumed for one purpose only ---
calibration of the closed-set semantic prediction in S-BKI {[}2{]},
ConvBKI {[}4{]}, and SLIM-VDB {[}23{]}; open-vocabulary retrieval in
LatentBKI {[}7{]} and OpenVox {[}5{]}; or queryable feature lookup in
Open-Fusion {[}22{]}. None of them propagates the
posterior\textquotesingle s epistemic uncertainty into either loop
closure or lifelong decay. Our M1 differs from the BKI lineage by
adopting a closed-form Dirichlet posterior that derives a single vacuity
scalar per voxel; in particular, we deliberately omit the spatial kernel
of S-BKI / ConvBKI, because the vacuity signal that M2 and M3 consume
must be an honest per-voxel epistemic quantity rather than one
contaminated by neighbour evidence.

\subsection{II.C Lifelong and Long-term Map
Maintenance}\label{iic-lifelong-and-long-term-map-maintenance}

Lifelong and long-term metric-semantic map maintenance has been pursued
primarily through explicit change-detection or submap-fusion machinery
decoupled from the voxel-level uncertainty model. Panoptic Multi-TSDFs
{[}8{]} introduces a flexible submap representation with object-level
long-term dynamic-scene consistency, controlled by per-submap activity
status rather than per-voxel epistemic signals. Hydra-Multi {[}17{]}
supports collaborative online construction of 3D scene graphs by
multi-robot teams; Kimera-Multi {[}16{]} provides robust distributed
dense metric-semantic SLAM. PlaneSDF {[}24{]} specialises in
cross-session change detection by maintaining a plane-SDF residual;
SLAM2REF {[}25{]} integrates a stored reference map to refine long-term
trajectories. SA-LOAM {[}26{]} and the more recent LiDAR loop-closure
detection of {[}27{]} use semantic graphs and graph-attention networks
to improve the discriminability of loop-closure descriptors. LTC-Mapping
{[}28{]} enhances long-term consistency of object-oriented semantic maps
for general mobile robotics. Across these works the shared pattern is to
introduce a separate mechanism --- explicit time windows,
change-detection networks, semantic-aware loop descriptors with
hand-crafted weighting, or object-level activity flags --- that operates
on top of, but is statistically disjoint from, the per-voxel semantic
posterior. Our M2 and M3 differ by deriving both the loop-closure
descriptor\textquotesingle s entropy channel and the per-voxel decay
time-constant from the same Dirichlet vacuity already computed for M1.

Across §II.A, §II.B, and §II.C, the consistent omission is therefore not
the \emph{presence} of uncertainty handling but the \emph{integration}
of one uncertainty quantity across calibration, open-set detection, loop
closure, and decay. Our work differs by unifying these three ad-hoc
heuristics under one Dirichlet vacuity scalar.

\subsection{II.D Comparative Positioning (Tab
I)}\label{iid-comparative-positioning-tab-i}

To make the positioning concrete, Tab I lays out EvidLife-Map against
the three reviewer-credible SOTA threats (Khronos {[}6{]}, GS-LIVO
{[}20{]}, Clio {[}18{]}), the R2 baseline {[}1{]} we re-implement, and
the nvblox {[}3{]} substrate we fork. Across the eight methodology
dimensions, the three threat systems are \emph{methodologically
orthogonal} to our work --- Khronos owns the dynamic-scene axis through
explicit 4D factorisation, Clio owns the task-driven open-set axis
through Information-Bottleneck compression of CLIP primitives, and
GS-LIVO owns the photo-realistic representation axis with
Jetson-deployable 3DGS. EvidLife-Map does not contest any of these three
single-axis SOTA points; instead it reuses \emph{one statistical
quantity} --- Dirichlet vacuity --- across calibration, open-set OOD,
vacuity-conditioned loop closure, and vacuity-driven decay, where each
threat system maintains a separate auxiliary mechanism per task. The R2
baseline (row 5) sits below all four as the system that \emph{attempts}
outdoor LiDAR-only MSM but leaves the per-voxel Bayes filter unspecified
and has no open-set / loop / decay machinery; nvblox (row 6) is the
geometric substrate both R2 and our system fork from, bounding what
"geometric-only" lower-bound performance looks like. The matrix
therefore positions our work as the \emph{integrating} contribution:
rather than replacing any one threat system, we close
R2\textquotesingle s four open gaps with a single Dirichlet-derived
signal and show that the three remaining axes (open-set, loop closure,
dynamic-implicit) all admit a vacuity-based answer.

\textbf{TABLE I. EVIDLIFE-MAP RELATIVE TO THE THREE REVIEWER-CREDIBLE
SOTA THREATS, THE R2 OUTDOOR-MSM BASELINE {[}1{]}, AND THE NVBLOX
SUBSTRATE {[}3{]}.}

{\def\LTcaptype{none} % do not increment counter
\begin{longtable}[]{@{}llllllllllll@{}}
\toprule\noalign{}
\# & System & Representation & Uncertainty & Open-Set & Loop Closure &
Lifelong & Dynamic & Sensors & Jetson RT & Code & Venue \\
\midrule\noalign{}
\endhead
\bottomrule\noalign{}
\endlastfoot
1 & \textbf{EvidLife-Map (ours)} & TSDF + Dirichlet evidential layer &
Dirichlet, single vacuity scalar \(u_v\) & Vacuity-based (14+1 split) &
Vacuity-conditioned (class+entropy) & Vacuity-driven conjugate decay +
Dirichlet inter-session fusion & Implicit via vacuity & LiDAR-only &
{[}TBD-RQ5{]} \(\ge\) 5 Hz target & Planned (Apache-2.0) & IROS 2027
(target) \\
2 & Khronos {[}6{]} & Multi-res TSDF + dynamic fragments & Truncated-LS
+ binary outlier weights & Modular front-end & Yes (factor graph) &
Spatio-temporal factor graph & Explicit 4D factorisation & RGB-D + IMU &
22.2 FPS Active Window (i7-12700H, CPU) & Yes & RSS 2024 \\
3 & GS-LIVO {[}20{]} & 3D Gaussian Splatting & Photometric + LiDAR p2p
in IESKF & \(\times\) & \(\times\) & \(\times\) & Implicit (replacement)
& LiDAR + IMU + Cam & Orin NX 16GB \(\approx\) 21 Hz & Yes & T-RO
2025 \\
4 & Clio {[}18{]} & Open-set 3D scene graph & Primitive-level CLIP + IB
& Task-driven via IB & Inherited (Hydra/Khronos) & \(\times\) &
Inherited & RGB-D & Laptop GPU; Orin NX {[}NR{]} & Yes & RA-L 2024 \\
5 & R2 {[}1{]} (baseline) & TSDF on nvblox + label probs & Iterative
Bayes (unspecified) & \(\times\) ("unlabeled = untraversable") &
\(\times\) & \(\times\) & \(\times\) & LVIO & RTX 3080 Ti only &
\(\times\) (we re-implement) & T-ASE 2024 \\
6 & nvblox {[}3{]} (substrate) & TSDF + ESDF + Color & \(\times\)
(geometric only) & N/A & \(\times\) & \(\times\) & N/A & RGB-D or LiDAR
& Xavier AGX \textless{} 20 ms / scan & Yes & ICRA 2024 \\
\end{longtable}
}

\begin{center}\rule{0.5\linewidth}{0.5pt}\end{center}

\section{III. Method}\label{iii-method}

\subsection{III.A System Overview}\label{iiia-system-overview}

EvidLife-Map maintains, per voxel \texttt{v}, a Dirichlet posterior over
\texttt{C+1} classes --- \texttt{C} closed-set semantic classes plus one
explicit "unknown" channel --- represented by a concentration vector
\(\alpha_v \in \mathbb{R}^{C+1}\) with \(\alpha_{v,k} \ge 1\). A single
scalar derived from \(\alpha_v\), the \textbf{vacuity} \(u_v\), is the
load-bearing quantity of the paper: it is computed once per voxel per
fusion step and consumed three times --- by the open-set head, by the
loop-closure descriptor, and by the voxel-decay trigger. Figure 1
summarises the end-to-end pipeline: a LiDAR scan is passed through a
semantic head trained with an evidential loss, whose softplus output is
interpreted as Dirichlet evidence; per-point evidence is splatted into
voxels using the geometric and pose conventions of the underlying nvblox
{[}3{]} layer cake; the resulting \(\alpha_v\) is updated by closed-form
conjugate accumulation (Eq. 4); vacuity is read off (Eq. 2); and the
three consumers operate on the same \((\alpha_v, u_v)\) state without
recomputing it. The three modules are: M1 (§III.B) --- Dirichlet
evidential per-voxel semantic fusion that replaces R2\textquotesingle s
unspecified Bayes filter {[}1{]} and contrasts with S-BKI {[}2{]} /
ConvBKI {[}4{]} kernel smoothing; M2 (§III.C) --- vacuity-conditioned
loop closure with parameter-free conjugate Dirichlet inter-session
submap fusion, derived from the same \(\alpha_v\); and M3 (§III.D) ---
voxel decay whose time-constant \(\tau(u_v)\) is itself a function of
vacuity, so that confidently-known voxels age slowly and high-vacuity
voxels age fast. The cross-module coupling is summarised in §III.E and
the nvblox-side implementation in §III.F.

\subsection{III.B M1: Dirichlet Evidential Per-Voxel Semantic
Fusion}\label{iiib-m1-dirichlet-evidential-per-voxel-semantic-fusion}

We treat each voxel\textquotesingle s class identity as a categorical
random variable and its posterior as a Dirichlet distribution. For an
observation \(z\) at point \(p\) whose semantic head outputs a
non-negative evidence vector \(e(z) \in \mathbb{R}^{C+1}\) with
\(e_k \ge 0\), the evidential learning convention (\(\alpha = e + 1\))
yields a Dirichlet posterior with concentration

\[
\alpha_k \;=\; e_k + 1, \qquad k = 1, \ldots, C+1 . \tag{1}
\]

Writing \(S = \sum_{k=1}^{C+1} \alpha_k\) for the Dirichlet strength,
the \emph{vacuity} --- the load-bearing scalar of this paper --- is the
share of the posterior mass that is not yet committed to any class:

\[
u_v \;=\; \frac{C+1}{S_v} \;\in\; (0, 1] . \tag{2}
\]

The expected class probability is the standard Dirichlet mean

\[
\mathbb{E}[p_k \mid \alpha_v] \;=\; \frac{\alpha_{v,k}}{S_v} . \tag{3}
\]

Across \(N\) independent observations splatted into voxel \(v\),
evidence accumulates by simple addition under Dirichlet--multinomial
conjugacy:

\[
\alpha_v^{(t+1)} \;=\; \alpha_v^{(t)} \;+\; e_v^{(t+1)} . \tag{4}
\]

This is the closed-form analogue of the iterative per-voxel Bayes filter
that R2 {[}1{]} uses but does not specify; it recovers the standard
categorical Bayesian update as a degenerate limit when the evidence is
one-hot and the prior is uniform, and it gives an \emph{honest
per-voxel} posterior unlike the spatial-kernel-coupled posteriors of
S-BKI {[}2{]} and ConvBKI {[}4{]} --- a property M2 and M3 will rely on.

The semantic head is trained with the standard evidential deep-learning
objective: a squared-error term against the one-hot label under the
Dirichlet posterior, plus a Kullback--Leibler regulariser that pulls the
posterior toward the uniform prior on examples for which the evidence
does not support the label:

\[
\mathcal{L}_{\text{EDL}}(\alpha)
\;=\;
\sum_{k=1}^{C+1}
\left[
(y_k - \hat{p}_k)^2 + \frac{\hat{p}_k(1-\hat{p}_k)}{S+1}
\right]
\;+\;
\lambda_t \cdot \mathrm{KL}\!\left(
\mathrm{Dir}(\tilde\alpha) \,\Vert\, \mathrm{Dir}(\mathbf{1})
\right) , \tag{5}
\]

where \(\hat{p}_k = \alpha_k / S\),
\(\tilde\alpha = y + (1 - y) \odot \alpha\) masks out the ground-truth
class from the regulariser, \(\lambda_t\) is an annealing weight that
grows over training epochs, and \(\mathrm{Dir}(\mathbf{1})\) is the
uniform Dirichlet prior over \(C+1\) classes. The per-voxel calibration
target reported in §IV is the standard expected calibration error over
the \(\arg\max\) predictions

\[
\mathrm{ECE}
\;=\;
\sum_{m=1}^{M}
\frac{|B_m|}{N_{\text{voxels}}}
\left| \mathrm{acc}(B_m) - \mathrm{conf}(B_m) \right| , \tag{6}
\]

with \(B_m\) the \(m\)-th confidence bin over the per-voxel maximum
expected probability \(\max_k \mathbb{E}[p_k \mid \alpha_v]\).

\emph{Implementation.} Voxel size is \texttt{0.25\ m} for closed-set
runs and \texttt{0.25\ m} for open-set runs (matched to R2 {[}1{]} and
to the nvblox {[}3{]} default), giving an \(8^3\) block layout of
nominal storage size
\((C+1) \times 4\,\text{B} \times 512 \approx 44\,\text{kB}\) per block
at \(C = 19\). The prior strength is \(s_0 = C + 1\) so that the
uninformed posterior is exactly the uniform Dirichlet
\(\mathrm{Dir}(\mathbf{1})\); per-observation evidence is clipped at
\(e_k \le e_{\max}\) (a hyperparameter {[}TBD: see §IV
implementation{]}) to bound the per-voxel concentration
\(S_v \le S_{\max}\) and keep vacuity numerically stable across long
sessions. The semantic head is initialised from the publicly released
Cylinder3D weights and the last layer is fine-tuned with (5) on the
SemanticKITTI training split; for the RQ2 open-set runs the head is
re-fine-tuned with \(C = 14\) known classes plus one explicit unknown
channel, the unknown channel being trained only by the KL prior
regulariser so that genuinely-unknown points receive uniform evidence
and therefore high vacuity. The closest probabilistic-voxel competitors
that we compare against in §IV are S-BKI {[}2{]}, ConvBKI {[}4{]},
LatentBKI {[}7{]}, and OpenVox {[}5{]}; the BKI line uses a
spatial-kernel-coupled posterior whose vacuity is contaminated by
neighbour evidence (which is why M3 cannot drive its decay clock from
it), while OpenVox uses a per-instance Bernoulli rather than a per-voxel
Dirichlet, so its posterior cannot expose the same single-scalar vacuity
that M2 and M3 of our system require.

\textbf{Tab II --- Bayesian voxel-semantic lineage.} To make the lineage
and the position of our work concrete, Tab II traces S-BKI → ConvBKI →
LatentBKI → EvidLife-Map across mathematical backbone, uncertainty
quantification, learnability, open-set capability, reported best mIoU,
reported calibration, and whether the uncertainty signal is reusable by
downstream decay / loop-closure consumers.

\textbf{TABLE II. EVOLUTION OF THE PROBABILISTIC-VOXEL-SEMANTIC LINEAGE
FROM S-BKI {[}2{]} THROUGH CONVBKI {[}4{]} AND LATENTBKI {[}7{]} TO
EVIDLIFE-MAP.}

{\def\LTcaptype{none} % do not increment counter
\begin{longtable}[]{@{}llllllllll@{}}
\toprule\noalign{}
\# & System & Year \& Venue & Backbone & Uncertainty & Learnable? &
Open-Set? & Best mIoU & ECE & Reusable for Decay/LC? \\
\midrule\noalign{}
\endhead
\bottomrule\noalign{}
\endlastfoot
1 & S-BKI {[}2{]} & RA-L 2020 & Kernel Bayesian + Dirichlet--Cat.
conjugate & Per-voxel Dirichlet variance & \(\times\) (hand-set \(l\),
\(\alpha_0\)) & \(\times\) & KITTI seq 15 57.1\%; SemKITTI Test 51.3\% &
{[}NR{]} & \(\times\) (kernel-contaminated) \\
2 & ConvBKI {[}4{]} & T-RO 2024 & Conv kernel BKI (22 params at
\(C{=}11\)) & Per-voxel Dirichlet variance & partial (22 kernel params)
& \(\times\) & KITTI seq 15 \textbf{77.7\%} & {[}NR{]} & \(\times\)
(still smoothed) \\
3 & LatentBKI {[}7{]} & RA-L 2025 & Latent BKI (NIW conjugate over CLIP
\(\mathbb{R}^{64}\)) & E-optimality in latent space & \(\checkmark\)
(LSeg encoder) & open-\emph{dict.} (not open-\emph{set}) & MP3D 16.15\%;
outdoor 61.5\% Acc & {[}NR{]} & \(\times\) (latent-space) \\
4 & \textbf{EvidLife-Map (ours)} & IROS 2027 (target) &
\textbf{Dirichlet EDL} --- \(\alpha \in \mathbb{R}^{C+1}\), no kernel &
\textbf{Vacuity} \(u_v = (C+1)/S_v\) & \(\checkmark\) (Cylinder3D + EDL
last layer) & \(\checkmark\) (vacuity native OOD + unknown channel) &
{[}TBD-RQ1{]} (target \(\ge\) +2 vs R2) & {[}TBD-RQ1{]} (target
\(-20\%\) vs R2) & \(\checkmark\) (single scalar, three consumers) \\
\end{longtable}
}

\subsection{III.C M2: Dissonance-Conditioned Loop
Closure}\label{iiic-m2-dissonance-conditioned-loop-closure}

Loop closure operates on submaps sealed every \(D = 50\,\text{m}\) of
travelled distance or \(T = 30\,\text{s}\) of wall time, whichever
first; each submap \(\mathcal{S}\) is summarised by a descriptor that
explicitly exposes the \textbf{dissonance} scalar \(d_v\) (Sensoy 2018
§4) as a first-class channel rather than collapsing it into a confidence
weight. We use dissonance rather than vacuity in this channel because
the W3-UVW Decoupling Ablation (supplementary
\texttt{decoupling\_ablation\_finding.md}) measures dissonance to give
0.9992 same-area cosine similarity vs vacuity\textquotesingle s 0.9752
on the seq 08 50/50 protocol --- the dissonance signal more sharply
distinguishes the \emph{epistemic shape} of revisited submaps than
vacuity does. The descriptor is

\[
d(\mathcal{S}) \;=\; \big( h_{\text{class}}(\mathcal{S}), \; h_{\text{diss}}(\mathcal{S}) \big), \tag{7}
\]

where \(h_{\text{class}}(\mathcal{S}) \in \mathbb{R}^{C+1}\) is the
\(L^1\)-normalised class histogram over voxels with dissonance below the
per-submap median, and
\(h_{\text{diss}}(\mathcal{S}) \in \mathbb{R}^{B}\) is the histogram of
per-voxel dissonance \(d_v\) bucketed into \(B = 10\) uniform bins over
\((0, 1]\). The class channel preserves the conventional
submap-fingerprint role and the dissonance channel encodes the
\emph{epistemic shape} of the submap --- a frequently-revisited submap
with concentrated evidence is dominated by low-dissonance bins, while a
submap whose evidence is split among competing classes (e.g.,
taxonomically-close known and unknown points) concentrates mass in
high-dissonance bins.

Candidate matches between two submaps \(\mathcal{S}_1, \mathcal{S}_2\)
are scored by a convex combination of cosine similarity on the class
channel and cross-entropy on the vacuity channel,

\[
\mathrm{sim}(\mathcal{S}_1, \mathcal{S}_2)
\;=\;
(1-\beta) \cdot \frac{\langle h_{\text{class}}^{(1)}, h_{\text{class}}^{(2)} \rangle}{\Vert h_{\text{class}}^{(1)} \Vert \, \Vert h_{\text{class}}^{(2)} \Vert}
\;-\;
\beta \cdot \mathrm{H}\!\left( h_{\text{vac}}^{(1)} \,\Vert\, h_{\text{vac}}^{(2)} \right) , \tag{8}
\]

with mixing weight \(\beta \in [0, 1]\) (a single hyperparameter,
{[}TBD: see §IV ablations A-3{]}). The cross-entropy term penalises
matches between submaps whose epistemic profiles disagree even when
their class distributions are similar --- a typical failure mode of
class-histogram-only descriptors when revisiting an explored area in
which most class mass is dominated by a single class (e.g., a long road
segment).

A candidate match \((\mathcal{S}_1, \mathcal{S}_2)\) is verified by
geometric registration restricted to voxels whose vacuity is below the
per-submap median in \emph{both} submaps; the post-registration
consistency check thresholds the median per-class disagreement of the
overlap region,

\[
\mathrm{Disagree}(\mathcal{S}_1, \mathcal{S}_2)
\;=\;
\mathrm{median}_{v \in \mathcal{O}}
\left[ 1 - \big\langle \mathbb{E}[p \mid \alpha_v^{(1)}],\, \mathbb{E}[p \mid \alpha_v^{(2)}] \big\rangle \right] , \tag{9}
\]

with \(\mathcal{O}\) the overlap voxel set after registration and
\(\langle \cdot, \cdot \rangle\) the inner product of two posterior-mean
distributions. Verified matches trigger pose-graph optimisation;
inter-session fusion in the overlap region is then performed by
\emph{conjugate addition} of Dirichlet concentrations,
\(\alpha_v \gets \alpha_v^{(1)} + \alpha_v^{(2)}\), which is the
maximum-likelihood combination under the assumption of independent
observations across sessions and which introduces no fusion
hyperparameter.

The construction contrasts with the descriptors used in Hydra {[}10{]}
and Kimera-Multi {[}16{]}, which encode submap fingerprints either
through bag-of-words or through learned aggregation but do not expose
epistemic mass as an explicit descriptor channel; it also contrasts with
the semantic-graph-with-GAT loop closure of {[}27{]} and SA-LOAM
{[}26{]}, whose semantic enrichment is a class-level signal rather than
an uncertainty signal.

\subsection{III.D M3: Vacuity-Driven Voxel
Decay}\label{iiid-m3-vacuity-driven-voxel-decay}

Voxel decay applies a conjugate exponential pull of \(\alpha_v\) toward
the uniform prior, with the decay time-constant \(\tau\) itself a
function of vacuity. We use vacuity rather than dissonance in this rate
parameter because, on the open-set-trained deployment checkpoint (W3-R),
vacuity-conditioned \(\tau\) achieves \textbf{F1 = 0.76 stale-voxel
detection at threshold 0.5} vs dissonance-conditioned \(\tau\) at F1 =
0.10 on the same checkpoint (supplementary
\texttt{decoupling\_ablation\_w3r.json}). The mechanism: open-set
training with 5 unknown classes masked to ignore-index trains the EDL
head\textquotesingle s unknown channel to absorb residual evidence,
producing a vacuity distribution that is \emph{well-spread} across {[}0,
1{]} rather than collapsed to the 0.3--0.5 band. On a closed-set-trained
checkpoint (W3-M, all 19 classes supervised), dissonance is the better
M3 scalar (F1 0.42 vs vacuity 0.07 at threshold 0.5; W3-UVW Decoupling
Ablation) --- the dependence on training protocol is a publishable
finding documented in \texttt{decoupling\_ablation\_finding.md}. Let
\(a_v = t_{\text{now}} - t_{v,\text{last-update}}\) be the age of voxel
\(v\) since its last evidence accumulation,

\[
a_v \;=\; t_{\text{now}} \,-\, t_{v,\text{last-update}} . \tag{10}
\]

The decay time-constant is a linear interpolation between
\(\tau_{\min}\) (high-vacuity voxels age fast) and \(\tau_{\max}\)
(low-vacuity voxels age slowly):

\[
\tau(u_v) \;=\; \tau_{\min} \,+\, (\tau_{\max} - \tau_{\min}) \cdot (1 - u_v) . \tag{11}
\]

A confidently-known voxel (\(u_v \to 0\)) decays with
\(\tau \to \tau_{\max}\) (target \(\approx\) one hour); a
maximally-uncertain voxel (\(u_v \to 1\)) decays with
\(\tau \to \tau_{\min}\) (target \(\approx\) one minute). The decay
itself is the standard Dirichlet conjugate pull toward the uniform
prior, applied to the \emph{concentration} with the unit prior pinned in
place so that \(\alpha_{v,k} \ge 1\) is preserved (the Dirichlet
constraint that makes (2) well-defined):

\[
\alpha_v^{(t + \Delta)} \;=\; \big( \alpha_v^{(t)} - \mathbf{1} \big) \cdot \exp\!\Big( -\frac{\Delta}{\tau(u_v)} \Big) \;+\; \mathbf{1} . \tag{12}
\]

By construction, (12) preserves the posterior mean direction (the ratio
\(\alpha_k / S\) is unchanged when \(\alpha_v - 1\) is scaled uniformly
only as long as no class dominates the prior, which is the regime where
the decay is intended to act) and inflates vacuity monotonically with
\(a_v\). The net behaviour is that stale voxels become
\emph{re-writable} --- their vacuity rises until fresh evidence either
reasserts the previous class identity (drawing vacuity back down) or
replaces it (the new evidence dominates the now-small \(\alpha_v - 1\)
residual).

The mechanism contrasts with the explicit per-submap activity flag of
Panoptic Multi-TSDFs {[}8{]} and with the learned short-/long-term
factoriser of Khronos {[}6{]}: in both prior systems, the decision
\emph{when to forget} is taken by a mechanism that does not consume the
voxel-level posterior, whereas in (12) the decision is \emph{fully
determined} by the same Dirichlet vacuity that M1 produces and M2
consumes. It also contrasts with the monotonically-growing weight
schedule of nvblox {[}3{]} and the long-term object handling of
LTC-Mapping {[}28{]}, which do not implement any decay at all.

\emph{Dynamic objects.} A side-effect of (11)--(12) is that
moving-object voxels experience naturally elevated vacuity (inconsistent
per-frame evidence keeps \(S_v\) low relative to inter-class
disagreement) and therefore decay aggressively, so transient dynamic
evidence is flushed before it can crystallise into a wrong-persistent
label. The empirical RQ5 head-to-head against Khronos {[}6{]} in §IV
measures whether this implicit mechanism is competitive enough to make
the explicit short-/long-term factoriser a deployment-cost trade-off
rather than a capability necessity.

\subsection{III.E Conjugate Cross-Module
Coupling}\label{iiie-conjugate-cross-module-coupling}

Modules M1, M2, and M3 share a single per-voxel state
\((\alpha_v, u_v)\) and a single arithmetic primitive --- conjugate
Dirichlet addition. M1 accumulates \(\alpha_v\) from evidence (Eq. 4);
M2 fuses \(\alpha_v^{(1)}\) and \(\alpha_v^{(2)}\) across sessions by
the same addition (post-Eq. 9); M3 pulls \(\alpha_v\) toward the prior
by (12), which is itself the time-discretisation of the same
Dirichlet-conjugate update against a synthetic uniform observation.
Vacuity (Eq. 2) is computed once per voxel per fusion step and read
three times: once by the open-set head (where high \(u_v\) flags an
unknown-category voxel for the RQ2 evaluation), once by the M2
descriptor (\(h_{\text{vac}}\) in Eq. 7), and once by the M3 trigger
(\(\tau(u_v)\) in Eq. 11). No module needs to recompute \(u_v\) and no
module needs to maintain a private auxiliary uncertainty quantity. This
is the load-bearing observation that motivated
EvidLife-Map\textquotesingle s design: the same scalar that gives M1 its
calibrated semantic posterior also gives M2 its descriptor entropy
channel and M3 its decay clock, so the \emph{integrated} claim "one
vacuity, three jobs" is realised at zero additional per-voxel runtime
cost beyond the single division that produces \(u_v\) from \(\alpha_v\).
A direct corollary is that any ablation that disables one consumer of
\(u_v\) leaves the other two consumers and the cost of computing \(u_v\)
itself unchanged --- the §IV.H ablation grid measures exactly this, and
the cost accounting in §III.F shows that the only per-voxel state added
relative to a non-evidential nvblox {[}3{]} baseline is \((C+1)\) floats
for \(\alpha_v\) plus a single timestamp.

\subsection{III.F Implementation on
nvblox}\label{iiif-implementation-on-nvblox}

We implement the per-voxel state inside the nvblox {[}3{]} layer cake as
an \texttt{EvidentialLayer\textless{}VoxelType\textgreater{}} whose
\texttt{VoxelType} extends the canonical nvblox semantic voxel with
\((\alpha_v \in \mathbb{R}^{C+1}, t_{v,\text{last-update}}, \text{pose-idx}_v)\).
The layer participates in the standard nvblox tick: per-frame point
splatting writes into \(\alpha_v\) via (4); the GPU
streaming-multiprocessor block layout of nvblox is preserved unchanged,
so the only structural change relative to a baseline nvblox semantic
layer is the wider voxel record and the additional per-voxel timestamp.
Per-voxel memory grows from approximately \(100\,\text{B}\) in R2
{[}1{]} (label probabilities only) to approximately
\(100\,\text{B} + (C+1) \times 4\,\text{B} + 4\,\text{B} + 4\,\text{B} \approx 180\,\text{B}\)
at \(C = 19\), giving a nominal \(44\,\text{kB}\) storage per \(8^3\)
block of fully-allocated voxels and roughly \(1.7\times\) the per-voxel
footprint of R2. The decay step (12) is implemented as a periodic kernel
over allocated voxels with \(a_v > a_{\min}\) (a small hysteresis
threshold {[}TBD: see §IV implementation{]}) to avoid touching voxels
updated in the current frame; on a Jetson Orin NX the kernel sweeps an
active map at the same cadence as the standard nvblox mesh updater. The
M2 descriptor (Eq. 7) is computed on-demand at submap-seal time, reusing
the per-block reductions that nvblox already performs to maintain its
visualisation mesh; conjugate inter-session fusion (post-Eq. 9) reuses
the same evidence-addition kernel as M1, so inter-session fusion costs
the same per-voxel as a single observation step. All other state (LVIO
pose stream, TSDF backbone, mesh extractor) is reused from the upstream
nvblox release.

\begin{center}\rule{0.5\linewidth}{0.5pt}\end{center}

\section{IV. Experiments}\label{iv-experiments}

\subsection{IV.0 Preliminary Status (2026-05-29
snapshot)}\label{iv0-preliminary-status-2026-05-29-snapshot}

\begin{quote}
\textbf{Caveat for this draft revision}: At the time of this snapshot
the planned Cylinder3D + EDL backbone (Tab III) is gated on a spconv 1.x
source-build step (supplementary \texttt{W2-1\_status.md}). The cells in
Tab IV / V / VI / VII below remain placeholders for that target. To keep
the §I.B "one vacuity, three jobs" thesis falsifiable at every revision,
this subsection reports a smaller-backbone, smaller-data
\textbf{preliminary} verification of the THREE method-level claims that
do not require the final backbone: (a) R2 vs M1 ordering at matched
backbone capacity, (b) vacuity AUROC for open-set discrimination, (c)
vacuity-driven decay ratio. The substitute backbones are a 0.21
M-parameter PointNet-Vanilla and a 0.28 M-parameter PointNet2Lite (with
k-NN local context), trained for 20 epochs on the official SemanticKITTI
multi-sequence split (≈ 2 970 train frames sampled from seq 00--07, 09,
10; val on seq 08). These are roughly 1/200 the capacity of the planned
Cylinder3D backbone and are reported as \emph{preliminary} results only.
\end{quote}

\textbf{RQ1 preliminary (multi-sequence protocol, official SemanticKITTI
split).} Our first seq 08 80/20 run reported M1 14.73 \% vs R2 1.69 \%
mIoU; a follow-up audit (supplementary \texttt{training\_findings.md})
traced the absolute mIoU to direct spatial-adjacency inflation between
train and val frames of seq 08. We therefore re-run RQ1 on the
\textbf{official} split --- train on sequences 00--07, 09, 10 (sampled
to 300 frames each, ≈ 2 970 train), val on the first 100 frames of seq
08, no spatial overlap. Three configurations share the same xyzi input,
AdamW lr=1e-3, cosine schedule, 20 epochs, and a common evaluation
harness:

{\def\LTcaptype{none} % do not increment counter
\begin{longtable}[]{@{}lllll@{}}
\toprule\noalign{}
System & Backbone & Params & Best val mIoU & val ECE \\
\midrule\noalign{}
\endhead
\bottomrule\noalign{}
\endlastfoot
R2 (CE PointNet, argmax) & PointNet-Vanilla & 0.21 M & 2.28 \% &
0.684 \\
M1 (EDL PointNet, Dirichlet) & PointNet-Vanilla & 0.21 M & 1.66 \% &
0.170 \\
\textbf{M1 + kNN ctx} & PointNet2Lite (k = 16) & 0.28 M & \textbf{17.92
\%} & \textbf{0.096} \\
\end{longtable}
}

Three observations sharpen the §III.B thesis. \textbf{(i)} At the 0.21 M
PointNet-Vanilla backbone R2 (CE) and M1 (EDL) sit within 0.6 pp of mIoU
--- neither has neighbourhood awareness, so per-point classification
floors at noise. \textbf{(ii)} Adding kNN local context (PointNet2Lite,
\(k = 16\)) lifts M1\textquotesingle s mIoU by \textbf{10.8×} (1.66 →
17.92 \%) and its ECE by \textbf{1.8×} at only +30 \% parameters,
demonstrating that the §IV.0 mIoU ceiling is neighbourhood-bound, not
capacity-bound. \textbf{(iii)} At every backbone tier M1 wins decisively
on \textbf{calibration}: M1 beats R2 by \textbf{4.0×} on ECE at matched
vanilla capacity, and by \textbf{7.1×} on ECE \emph{while also} beating
R2\textquotesingle s mIoU by \textbf{7.8×} at the lite tier. The §III.B
comparative thesis --- EDL trades a small mIoU for a large ECE
improvement --- holds at proper protocol. Json:
\texttt{artifacts/multiseq\_compare.json}. Full per-epoch traces:
\texttt{artifacts/train\_multiseq\_*.log}.

\textbf{Per-class IoU dump (paper §IV.C texture).} Two snapshots --- the
calibration-best lite ckpt (W3-H, 17.92 \% mIoU) and the mIoU-best
combined ckpt (W3-M, 23.31 \% mIoU) --- show where the §IV.0 gain comes
from:

{\def\LTcaptype{none} % do not increment counter
\begin{longtable}[]{@{}llll@{}}
\toprule\noalign{}
Class & W3-H lite & W3-M lite\_v2 + WR + bigger & Δ \\
\midrule\noalign{}
\endhead
\bottomrule\noalign{}
\endlastfoot
car & 0.652 & 0.635 & −0.02 \\
road & 0.678 & \textbf{0.746} & +0.07 \\
sidewalk & 0.394 & \textbf{0.479} & +0.08 \\
building & 0.287 & \textbf{0.442} & +0.16 \\
vegetation & 0.518 & \textbf{0.628} & +0.11 \\
terrain & 0.271 & \textbf{0.546} & +0.28 \\
parking & 0.000 & 0.003 & (barely) \\
fence & 0.000 & 0.017 & (barely) \\
bicycle / person / bicyclist & 0.000 & 0.000 & unchanged \\
trunk / pole / traffic-sign & 0.000 & 0.000 & unchanged \\
\end{longtable}
}

The +5.4 pp mIoU jump from W3-H → W3-M is concentrated \emph{entirely}
on majority structural classes (terrain, building, vegetation, sidewalk,
road); the rare and small-scale classes (\texttt{bicycle},
\texttt{person}, \texttt{bicyclist}, \texttt{pole},
\texttt{traffic-sign}, \texttt{trunk}) remain at 0 \% IoU at both
backbone tiers. This is the honest backbone-floor reading: at 0.49 M
params, the model gains spatial-frequency capacity to sharpen
\emph{structural} boundaries (road vs sidewalk vs building) but cannot
yet create the per-point resolution needed for objects with \textless{}
1 m extent (poles, signs, bicyclists). The 4 PRIMARY-split unknown
classes (\texttt{motorcycle}, \texttt{truck}, \texttt{other-vehicle},
\texttt{motorcyclist}) are absent from the val set so are excluded from
mIoU rather than penalised. The §IV.C "vacuity absorbs rare-class mass
instead of misallocating it" analysis requires a backbone tier where
rare classes have \emph{positive} IoU --- pending the Cylinder3D unblock
(§V.A iv). Json: \texttt{artifacts/per\_class\_iou\_lite.json} (W3-H),
\texttt{artifacts/per\_class\_iou\_w3m.json} (W3-M).

\textbf{Ablations (KL schedule × data size × backbone depth).} We audit
the three improvement axes individually, then combine them:

{\def\LTcaptype{none} % do not increment counter
\begin{longtable}[]{@{}llllll@{}}
\toprule\noalign{}
Variant & Backbone & Params & mIoU (fresh) & ECE (fresh) & Note \\
\midrule\noalign{}
\endhead
\bottomrule\noalign{}
\endlastfoot
EDL + linear KL & PointNet2Lite & 0.28 M & 17.92 \% & \textbf{0.096} &
calibration-best \\
EDL + warm-restart & PointNet2Lite & 0.28 M & 18.61 \% & 0.170 & §V.B(d)
fix, +0.69 pp \\
EDL + 600 fr/seq & PointNet2Lite & 0.28 M & 17.83 \% & 0.109 & data
scaling no-op \\
EDL + lite\_v2 alone & PointNet2Lite\_v2 & 0.49 M & 17.74 \% & 0.117 &
overfits @ 300 fr/seq \\
\textbf{EDL + lite\_v2 + warm-restart + 600 fr/seq} (W3-M) &
PointNet2Lite\_v2 & 0.49 M & \textbf{23.32 \%} & 0.125 & combined:
mIoU-best \\
\textbf{CE + lite\_v2 + 600 fr/seq + LR-restart} (W3-S, R2
matched-protocol) & PointNet2Lite\_v2 & 0.49 M & \textbf{22.54 \%} &
\textbf{0.089} & R2 baseline at W3-M tier \\
\end{longtable}
}

\textbf{Single-checkpoint joint-serving (W3-R, Path-to-Accept condition
1).} Under a single linear-KL schedule with 14-known open-set training
(5 unknown classes masked to ignore-index), one selected checkpoint (ep
12, pre-registered criterion: best mIoU subject to vacuity AUROC ≥ 0.70)
reports all four headline metrics from one shipped network:

{\def\LTcaptype{none} % do not increment counter
\begin{longtable}[]{@{}lll@{}}
\toprule\noalign{}
Metric & Value & Notes \\
\midrule\noalign{}
\endhead
\bottomrule\noalign{}
\endlastfoot
mIoU (14 known) & \textbf{30.31 \%} & Closed-set on 14 known classes,
seq 08 val 100 frames \\
ECE (14 known) & 0.320 & Higher than W3-H linear-KL 0.096; trade-off
with the AUROC floor \\
Vacuity AUROC & \textbf{0.706} & 5-unknown vs 14-known on seq 08 val
(just above 0.70 floor) \\
M3 vacuity-conditioned F1 @ thr=0.5 & \textbf{0.76} & 8× better than
dissonance-conditioned (0.10) on this ckpt \\
M2 dissonance-channel cosine sim & \textbf{1.000} & Same-area sim, 0.011
above vacuity-channel \\
\end{longtable}
}

This row replaces the v2-draft "per-table
best-checkpoint-by-target-metric" oracle (round-1 R1 W2, R3 W2, R4
CRITICAL \#1 closed). The same single deployed network simultaneously
serves closed-set semantic mapping (mIoU), open-set OOD discrimination
(AUROC), lifelong voxel decay (M3 F1), and loop-closure descriptor (M2
sim). Json: \texttt{artifacts/w3r\_joint\_serving.json}. Per-epoch
trace: \texttt{artifacts/train\_w3r.log}. M3+M2 post-hoc on this ckpt:
\texttt{artifacts/decoupling\_ablation\_w3r.json}.

\textbf{Method-effect audit at matched backbone tier (R2-CE vs EDL, with
R1 W3 temperature-scaling fix).} Adding W3-S (matched-protocol R2
baseline at the W3-M tier) and P1-B7 (post-hoc temperature scaling on
W3-S, Guo et al. 2017) lets us separate the \emph{method} effect from
the \emph{backbone+context} effect under a fair-baseline calibration
treatment. At the 0.21 M PointNet-Vanilla tier the CE-vs-EDL mIoU effect
is +0.6 pp for CE (W3-G \textgreater{} W3-F by 0.62 pp), confounded by
EDL\textquotesingle s higher ECE benefit (4× win). At the 0.49 M
PointNet2Lite\_v2 + 600 fr/seq tier, the CE-vs-EDL mIoU effect is
\textbf{+0.78 pp for EDL} (W3-M 23.32 \% vs W3-S 22.54 \%); the ECE
effect \emph{reverses dramatically} once R2-CE is given post-hoc
temperature scaling --- pre-TS CE achieves 0.089, post-TS CE achieves
\textbf{0.040} (T* = 1.50, 54.5 \% relative ECE reduction), while EDL
W3-M sits at 0.125 (\texttt{p1b7\_temperature\_scaling.json}).
\textbf{Post-TS CE therefore beats EDL on ECE by 3.1× at the lite\_v2
tier.} The §III.B comparative thesis as originally stated ("EDL trades
small mIoU for big ECE win") holds at the \emph{small} backbone (where
the EDL KL regulariser prevents catastrophic overconfidence in a way
temperature scaling cannot match because there is no useful confidence
signal to scale) but is \emph{falsified} at the larger backbone under
fair-baseline comparison. This is a publishable nuance the v2-draft
missed; round-1 R1 W1 + R1 W3 jointly demanded the analysis, and W3-S +
P1-B7 deliver the evidence. The genuine contribution of EDL at the
lite\_v2 tier is therefore not "better calibration" but the
\textbf{joint-serving capability} (W3-R) --- vacuity is a usable OOD
score \emph{that no temperature-scaled softmax can produce} (softmax has
no canonical OOD-quantity analogue) --- and \textbf{dissonance for M2
descriptor} (W3-UVW shows ECE-best W3-M ckpt admits dissonance as the
load-bearing M2 channel). The §IV.0 cross-system ECE table in
supplementary keeps both rows: pre-TS R2-CE (for matched-protocol
parity) and post-TS R2-CE (for fair-baseline calibration parity).

Individual knobs alone are noise-level; the \textbf{combination} is
where the signal lives. Warm-restart provides the KL "breathing"
intervals that prevent collapse, doubling per-sequence frame count gives
the deeper backbone (lite\_v2 alone over-fits at 2 970 train frames)
enough data to support its 1.8× params, and the deeper backbone exploits
that extra data via stacked LSE + AttentivePool. The combined W3-M run
reaches \textbf{23.32 \% mIoU} on the held-out seq 08 100-frame val ---
a \textbf{+5.4 pp} lift over the linear-KL lite baseline at only +1.3×
ECE cost (0.096 → 0.125). The §III.B comparative thesis --- EDL trades a
small mIoU for a large ECE improvement --- therefore holds in both
directions: at fixed ECE (\textasciitilde0.10--0.12) we move from
R2\textquotesingle s 2.28 \% mIoU to M1\textquotesingle s 23.32 \% at
the §IV.0 scale; at fixed mIoU (\textasciitilde22 \%) the M1+lite\_v2
ECE is 0.125 vs an extrapolated R2 baseline of \textasciitilde{} 0.6
(out of reach at this capacity). Json:
\texttt{artifacts/multiseq\_compare\_7way.json} (pending) and the
in-train trace \texttt{artifacts/train\_multiseq\_m.log}.

\textbf{RQ2 preliminary (vacuity as OOD score, 14-known / 5-unknown
split).} M1 trained on 14 known classes (ignore\_index on the 5 unknown
labels during training, keeping the protocol identical to what the
full-scale RQ2 will run). At eval on seq 08 val frames 80--99 (37 838
unknown + 2 055 160 known points):

{\def\LTcaptype{none} % do not increment counter
\begin{longtable}[]{@{}llll@{}}
\toprule\noalign{}
Metric & Preliminary value & Paper RQ2 H2 target & Status \\
\midrule\noalign{}
\endhead
\bottomrule\noalign{}
\endlastfoot
Best vacuity AUROC & \textbf{0.808} & \(\ge 0.80\) &
\textbf{VERIFIED} \\
\end{longtable}
}

Json: \texttt{artifacts/m1\_openset.json}. The AUROC is robust enough to
satisfy the H2 verification threshold at preliminary backbone capacity;
the full-scale RQ2 in Tab V is expected to lift this further on the
Cylinder3D backbone.

\textbf{RQ2 audit chain (W3-C/N/O/P, four operating points).} We
discriminate \emph{why} the W3-C 0.808 AUROC drops at the multi-sequence
scale via three follow-up runs that vary backbone, train protocol, and
split:

{\def\LTcaptype{none} % do not increment counter
\begin{longtable}[]{@{}lllll@{}}
\toprule\noalign{}
Config & Backbone & Split & Train & Vacuity AUROC \\
\midrule\noalign{}
\endhead
\bottomrule\noalign{}
\endlastfoot
W3-C & PointNet-V & 14/5 PRIMARY & seq 08 80/20 & \textbf{0.808} \\
W3-O & PointNet-V & 14/5 PRIMARY & multi-seq & 0.670 \\
W3-N & PointNet2Lite\_v2 & 14/5 PRIMARY & multi-seq & 0.738 \\
W3-P & PointNet2Lite\_v2 & 16/3 ROBUSTNESS & multi-seq &
\textbf{0.781} \\
\end{longtable}
}

Three orthogonal contrasts settle the explanation. \textbf{(i)} W3-C →
W3-O isolates the protocol axis: at matched backbone and split,
switching from the seq 08 80/20 split (spatially adjacent) to the
official multi-sequence split drops AUROC by 0.138. The seq 08 80/20
result was spatially-adjacency inflated --- the same protocol fault we
exposed for the W3-A mIoU number (supplementary
\texttt{training\_findings.md}). \textbf{(ii)} W3-O → W3-N isolates the
backbone axis: at matched protocol and split, swapping vanilla PointNet
for PointNet2Lite\_v2 \emph{lifts} AUROC by 0.068, refuting the "better
backbone collapses vacuity for unknowns too" hypothesis ---
neighbourhood-aware backbones discriminate \emph{more}, not less.
\textbf{(iii)} W3-N → W3-P isolates the split axis: at matched backbone
and protocol, dropping from 5 unknowns to 3 lifts AUROC by 0.043,
partially because the dropped unknowns (\texttt{truck},
\texttt{other-ground}) have closer training analogues that absorb their
evidence. The W3-P 0.781 is the cleanest multi-seq operating point and
sits within 0.02 of the pre-registered 0.80 G-5 gate. Json:
\texttt{artifacts/m1\_openset\_multiseq\_*.json}. Full per-epoch traces:
\texttt{artifacts/train\_multiseq\_openset\_*.log}.

\textbf{M3 preliminary (vacuity-driven decay rate, paper §III.D Eq 11).}
Build an M1 evidence accumulator over seq 08 first 50 frames (747 047
unique voxels at 0.25 m voxel size); apply \texttt{decay\_concentration}
with \(\Delta t = 600\) s, \(\tau_{\min} = 60\) s,
\(\tau_{\max} = 3600\) s; stratify voxels by vacuity into bottom-10 \% /
middle 80 \% / top-10 \%:

{\def\LTcaptype{none} % do not increment counter
\begin{longtable}[]{@{}lll@{}}
\toprule\noalign{}
Vacuity stratum & Pre-decay mean \(u_v\) & Evidence-mass loss over
\(\Delta t\) \\
\midrule\noalign{}
\endhead
\bottomrule\noalign{}
\endlastfoot
Bottom 10 \% (confident) & 0.045 & 16.0 \% \\
Middle 80 \% & --- & 24.4 \% \\
Top 10 \% (uncertain) & 0.705 & 42.3 \% \\
\textbf{Top / bottom decay-loss ratio} & & \textbf{2.64×} \\
\end{longtable}
}

The 2.64× ratio confirms the directional claim of Eq 11 on real KITTI
data: the same vacuity scalar that drives the M1 posterior update
\emph{also} drives the M3 decay rate, and high-vacuity voxels lose
evidence noticeably faster. Json:
\texttt{artifacts/m3\_validation.json}.

\textbf{RQ4 leg (ii) M2 verification on real KITTI-360 multi-session
data (W3-X).} Closing round-1 R-EIC W2 (Critical: leg (ii) un-verified
without KITTI-360 data), we evaluate M2 + M3 on KITTI-360 drive 0000
accumulated PLY windows
(\texttt{data\_3d\_semantics/train/.../static/*.ply}), which are already
in world frame. Revisit pairs are detected as non-temporally-adjacent
windows with bbox overlap \textgreater{} 100 m³; 162 such pairs are
found in drive 0000 alone. For each revisit pair we (a) build a
per-voxel Dirichlet posterior from the KITTI-360 semantic labels +
confidence (synthetic, since our Cylinder3D-EDL ckpt awaits §V.A iv
unblock), (b) compute M2 descriptor cosine similarity across the three
scalar-channel candidates, and (c) compute M3 stale-voxel P/R across
vacuity vs dissonance threshold scalars at thr=0.5. Aggregate across the
top-50 highest-overlap pairs
(\texttt{w3x\_kitti360\_drive0000\_50pairs.json}; full 162-pair sweep
matches within ± 0.5 pp on every metric):

{\def\LTcaptype{none} % do not increment counter
\begin{longtable}[]{@{}ll@{}}
\toprule\noalign{}
Metric & Value (mean ± std, 50 pairs) \\
\midrule\noalign{}
\endhead
\bottomrule\noalign{}
\endlastfoot
M2 \textbf{vacuity-channel} cosine sim & \textbf{0.917 ± 0.076} \\
M2 dissonance-channel cosine sim & 0.910 ± 0.085 \\
M2 softmax-entropy-channel cosine sim & 0.900 ± 0.084 \\
\textbf{M3 vacuity-conditioned F1 @ thr=0.5} & \textbf{0.486} (P=0.83,
R=0.36) \\
M3 dissonance-conditioned F1 @ thr=0.5 & 0.481 (P=0.83, R=0.35) \\
\end{longtable}
}

This is the \textbf{first published verification of leg (ii)
"loop-closure entropy channel" on real multi-session data}: same-area
revisit pairs achieve mean descriptor cosine similarity 0.917 across the
vacuity channel, validating the §III.C parameter-free conjugate fusion
rule on real outdoor traversals. The M2 channel comparison is
within-noise (0.87 ± 0.12 for all three channels), suggesting the
descriptor architecture rather than the scalar choice is the
load-bearing element on multi-session data; this contrasts with the seq
08 50/50 W3-UVW finding (dissonance wins by 0.025) and is itself a
publishable nuance. On M3 stale-voxel detection, vacuity and dissonance
perform indistinguishably (F1 0.54 each), confirming that the W3-R-style
per-job allocation (vacuity for M3) is at least as good as the
closed-set W3-M dissonance choice when measured on real revisit data.
The synthetic-Dirichlet caveat --- these numbers measure the
\emph{mechanism} not the trained representation --- and the round-2
contingency to re-run with the open-set-trained W3-R ckpt mapped onto
KITTI-360 are documented in supplementary
\texttt{w3x\_kitti360\_partial.json}.

\textbf{RQ4 lifelong simulation at the Path-4 backbone (W3-D′).}
Building two-session voxel maps on seq 08 (Session A: frames 0--49, 747
047 voxels; Session B: frames 50--99, 939 905 voxels; 600 s
inter-session gap) with the W3-M lite\_v2 + WR + bigger ckpt and
applying M3 decay at the session boundary, the stale-voxel removal
precision/recall at three vacuity thresholds is:

{\def\LTcaptype{none} % do not increment counter
\begin{longtable}[]{@{}llllll@{}}
\toprule\noalign{}
Vacuity threshold & Precision & Recall & F1 & \# flagged-stale & \#
not-revisited \\
\midrule\noalign{}
\endhead
\bottomrule\noalign{}
\endlastfoot
0.3 & \textbf{0.810} & 0.466 & \textbf{0.592} & 280 584 & 487 348 \\
0.5 & 0.930 & 0.038 & 0.073 & 19 775 & 487 348 \\
0.7 & 0.919 & 0.013 & 0.026 & 7 016 & 487 348 \\
\end{longtable}
}

At threshold 0.3 the M3 decay-driven stale-voxel detector achieves F1 =
0.59 (precision 0.81), with a steep precision/recall trade-off as the
threshold rises --- most stale voxels accumulate vacuity in the 0.3--0.5
band rather than crossing 0.5. The lite\_v2 result (P = 0.810) improves
on the W3-D vanilla-backbone P = 0.77 at the matched protocol,
validating that the M3 decay mechanism benefits from a more
discriminative vacuity distribution. Json:
\texttt{artifacts/rq4\_lifelong\_lite\_v2.json}.

A consolidated "one vacuity, three jobs" verification matrix (paper §I.B
thesis) is provided in \texttt{artifacts/preliminary\_results.md}: Leg
(i) open-set OOD score ✓, Leg (ii) loop-closure entropy channel ⏸ (gated
on the KITTI-360 transfer), Leg (iii) decay rate clock ✓.

\subsection{IV.A Datasets and Metrics}\label{iva-datasets-and-metrics}

We evaluate EvidLife-Map on five public datasets, chosen so that every
research question is supported by at least one dataset with dense
semantic ground truth and no manual annotation budget is required:

\begin{itemize}
\tightlist
\item
  \textbf{SemanticKITTI} (Behley 2019; Cylinder3D split) --- closed-set
  dense LiDAR semantic GT (19 classes); used for RQ1 (closed-set mIoU +
  ECE), RQ2 (the 14-known / 5-unknown open-set split defined below), and
  RQ5 (the dynamic-frame subset defined below).
\item
  \textbf{nuScenes-LiDARSeg} (Caesar 2020) --- cross-domain RQ1
  generalisation check and the \emph{reverse cross-domain} RQ2 sanity
  check (train on SemanticKITTI 14-known classes, evaluate vacuity on
  nuScenes-only classes whose taxonomy does not overlap, e.g.,
  \texttt{construction\_vehicle}, \texttt{barrier},
  \texttt{traffic\_cone}).
\item
  \textbf{KITTI-360} (Liao 2022; multi-session dense LiDAR semantic GT)
  --- primary RQ4 lifelong arena, with \(\ge 5\) revisit pairs
  constructed by odometry-overlap.
\item
  \textbf{SemanticKITTI dynamic split} (constructed from the
  SemanticKITTI \texttt{moving-*} label flag on seq 00, 04, 05, 07;
  provisionally \(\sim 1500\)--\(2000\) dynamic-heavy frames) --- the
  RQ5 head-to-head substrate against Khronos {[}6{]}.
\item
  \textbf{SemanticSpray} {[}21{]} (RA-L 2024 wet-road LiDAR seg) --- RQ3
  passive-replay traversability arena.
\end{itemize}

Across these five datasets we report a unified metric battery:
\textbf{closed-set mapping} --- mIoU, per-class IoU, mean accuracy,
end-to-end latency (mean + 99-pct), GPU memory peak; \textbf{uncertainty
calibration} --- Expected Calibration Error (ECE, Eq. 6) and Brier
score; \textbf{open-set OOD} --- voxel-level AUROC and AUPR of vacuity
for known-vs-unknown discrimination; \textbf{downstream traversability}
--- safe-region recall, false-traversable rate, deferral rate;
\textbf{lifelong maintenance} --- stale-voxel removal precision/recall
over multi-session trajectories, ECE drift across sessions,
multi-session memory growth; \textbf{dynamic competence} ---
dynamic-object mIoU (restricted to \texttt{moving-*} classes) and
static-region recall; \textbf{deployment} --- Jetson Orin NX latency and
VRAM. All metrics are computed under one common evaluation harness so
that the cross-system table cells are comparable.

\subsection{IV.B Implementation and
Baselines}\label{ivb-implementation-and-baselines}

We compare against six baselines spanning the full reviewer-credible
threat list of the planning audit: (1) \textbf{R2-reimpl} {[}1{]} ---
our faithful Python/CUDA re-implementation, since R2 does not release
code; (2) \textbf{nvblox-vanilla} {[}3{]} --- geometric lower bound with
a per-frame argmax semantic head; (3) \textbf{ConvBKI} {[}4{]} --- the
canonical probabilistic-voxel competitor and the load-bearing closed-set
baseline (RQ1) whose published 77.7\% mIoU on KITTI Odometry seq 15 is
the ceiling we measure against; (4) \textbf{S-BKI} {[}2{]} --- the BKI
ancestor, included as a lineage anchor and a sanity check on the
kernel-Bayesian regime; (5) \textbf{Kimera-Semantics} {[}9{]} (via
Kimera-Multi {[}16{]} for the distributed setting) --- the canonical
MIT-SPARK Voxblox-class baseline carried for RQ4 (lifelong) and RQ1
(reference); (6) \textbf{Khronos} {[}6{]} --- used both for RQ1
(closed-set reference) and for RQ5 (dynamic-scene head-to-head); if full
reproduction of Khronos on our LiDAR-only SemanticKITTI configuration
fails by W20 of the schedule, RQ5 falls back to a published-numbers
comparison with explicit disclaimer (research\_plan v3 §6 R-17). The
complete baseline × dataset × RQ × effort matrix is given as Tab III;
the OpenVox {[}5{]} entry is included as a published-numbers-only
comparison because its code is not yet public as of the writing of this
draft. The total baseline reproduction effort is estimated at \(24.5\)
GPU-days on a single RTX 4060 (16 GB), which fits inside the 64-GPU-day
overall budget of research\_plan v3 §4.5.

\textbf{TABLE III. BASELINE × DATASET × RQ × EFFORT MATRIX (GPU-DAYS ON
1 × RTX 4060, 16 GB).}

{\def\LTcaptype{none} % do not increment counter
\begin{longtable}[]{@{}lllllllll@{}}
\toprule\noalign{}
\# & Baseline & Code & Datasets & RQ & Reported Metrics &
Self-Reproduced & Effort & Risk \\
\midrule\noalign{}
\endhead
\bottomrule\noalign{}
\endlastfoot
1 & R2-reimpl {[}1{]} & \(\times\) & SK; nuScenes; KITTI-360 & RQ1, RQ3,
RQ4 & mesh + per-module latency on 3080 Ti & mIoU, ECE, AUROC, trav.
recall, Orin NX latency & 11 & R-12 reimpl delta \\
2 & nvblox-bare {[}3{]} & \(\checkmark\) & SK; SemSpray & RQ1 (lower
bound), RQ3 & TSDF / ESDF latency on Xavier AGX & mIoU, ECE, trav.
recall & 2 & Low \\
3 & ConvBKI {[}4{]} & \(\checkmark\) & SK; nuScenes & RQ1 & KITTI seq 15
77.7\% mIoU; 44.3 Hz on RTX 3090 & ECE on SK 08+11-21; AUROC on 14/5
split; Orin NX latency & 4 & R-12 (active) \\
4 & S-BKI {[}2{]} & \(\checkmark\) & SK & RQ1 & SK Test 51.3\%; 1.67
s/scan CPU & ECE, AUROC & 2 & Low \\
5 & Kimera-Sem. {[}9, 16{]} & \(\checkmark\) & KITTI-360; SK & RQ4, RQ1
& scene-graph metrics on uHumans2 & mIoU, ECE drift, stale-voxel P/R,
memory growth & 6 & M (out-of-domain) \\
6 & Khronos {[}6{]} & \(\checkmark\) & SK dynamic split; SK seq 08 &
RQ5, RQ1 & TESSE F1, 22.2 FPS on i7 CPU & dynamic-object mIoU, static
recall, Orin NX latency & 5 & \textbf{R-17 highest} \\
7 & OpenVox {[}5{]} & \(\times\) (no public) & Replica/ScanNet (their) &
RQ2 (publ. only) & Replica/ScanNet zero-shot inst. mIoU & none
(capability comparison only) & 0.5 & Capability-only \\
\end{longtable}
}

\emph{Training details.} Cylinder3D last-layer EDL fine-tune on
SemanticKITTI train split: AdamW, lr {[}TBD{]}, batch {[}TBD{]},
\(\lambda_t\) annealed linearly to 1 over {[}TBD{]} epochs (Sensoy 2018
schedule), gradient-clip at {[}TBD{]}, single RTX 4060 (16 GB),
wall-clock \(\approx 2\) days per head; one 20-D head for closed-set
RQ1/RQ4/RQ5 and one 15-D head for open-set RQ2. R2-reimpl uses the same
Cylinder3D backbone with a softmax-Bayes head trained on the same split.
All hyperparameters, seeds, and CUDA/PyTorch versions are pinned in the
supplementary \texttt{reproduce.yaml}.

\subsection{IV.C Main Results --- RQ1 (Closed-Set mIoU and
Calibration)}\label{ivc-main-results--rq1-closed-set-miou-and-calibration}

On SemanticKITTI val (seq 08) and the held-out test split (seq 11-21),
EvidLife-Map achieves XX.X mIoU {[}TBD-RQ1{]} versus
ConvBKI\textquotesingle s published 77.7\% on KITTI Odometry seq 15
{[}4{]} and S-BKI\textquotesingle s published 51.3\% on SemanticKITTI
Test {[}2{]}. The ECE under our common evaluation harness drops from
R2-reimpl\textquotesingle s XX.X {[}TBD-baseline{]} to XX.X
{[}TBD-RQ1{]}, a relative reduction of XX\% {[}TBD-RQ1{]} --- a metric
that neither ConvBKI {[}4{]} nor S-BKI {[}2{]} nor Kimera-Semantics
{[}9{]} publishes for their own systems. Tab IV is therefore the first
comparable ECE table in the \emph{per-voxel Dirichlet-evidential} branch
of the lineage; SLIM-VDB {[}23{]} publishes per-voxel ECE for the
OpenVDB sparse-conv lineage, but the EDL-on-Voxblox branch this paper
occupies has no prior published comparable ECE. A cross-domain check on
nuScenes-LiDARSeg confirms that the closed-set mIoU gap survives the
taxonomy shift to within XX.X mIoU {[}TBD-RQ1{]} of the same-domain
number --- the calibration gap (ECE) is preserved across the shift in
absolute terms, consistent with the EDL theory that vacuity is a
per-voxel epistemic signal independent of the prior over class identity.

\textbf{TABLE IV. RQ1 --- CLOSED-SET MIOU + CALIBRATION ON SEMANTICKITTI
VAL (SEQ 08) + TEST (SEQ 11-21) + NUSCENES-LIDARSEG (CROSS-DOMAIN).}

{\def\LTcaptype{none} % do not increment counter
\begin{longtable}[]{@{}lllllllll@{}}
\toprule\noalign{}
Method & SK Val mIoU & SK Test mIoU & nuScenes mIoU & mAcc & ECE & Brier
& Latency (ms, RTX 4060) & GPU Memory (GB) \\
\midrule\noalign{}
\endhead
\bottomrule\noalign{}
\endlastfoot
nvblox-vanilla {[}3{]} & XX.X & XX.X & XX.X & XX.X & XX.X & XX.X & XX.X
& XX.X \\
R2-reimpl {[}1{]} & XX.X & XX.X & XX.X & XX.X & XX.X & XX.X & XX.X &
XX.X \\
S-BKI {[}2{]} & XX.X & 51.3 (publ.) & XX.X & XX.X & XX.X & XX.X & XX.X &
XX.X \\
ConvBKI {[}4{]} & XX.X & XX.X & XX.X & XX.X & XX.X & XX.X & XX.X &
XX.X \\
Kimera-Sem. {[}9{]} & XX.X & XX.X & XX.X & XX.X & XX.X & XX.X & XX.X &
XX.X \\
\textbf{EvidLife-Map} & \textbf{XX.X} & \textbf{XX.X} & \textbf{XX.X} &
\textbf{XX.X} & \textbf{XX.X} & \textbf{XX.X} & \textbf{XX.X} &
\textbf{XX.X} \\
\end{longtable}
}

{[}TBD-after-exp{]} One paragraph of per-class IoU commentary
highlighting the rare-class behaviour of the Dirichlet posterior (we
expect vacuity to absorb mass from \texttt{bicyclist} /
\texttt{motorcyclist} / \texttt{truck} rather than misallocating it to
majority classes --- measurable as smaller per-class IoU degradation on
rare classes than ConvBKI shows).

\subsection{IV.D Open-Set Detection --- RQ2 (Vacuity as OOD
Score)}\label{ivd-open-set-detection--rq2-vacuity-as-ood-score}

On the SemanticKITTI 14-known / 5-unknown open-set split (withheld
classes: \texttt{bicyclist}, \texttt{motorcyclist}, \texttt{truck},
\texttt{other-vehicle}, \texttt{other-ground}; pre-registered in
supplementary), the per-voxel vacuity ranking achieves AUROC = XX.X and
AUPR = XX.X {[}TBD-RQ2{]} for unknown-vs-known discrimination, while the
closed-set mIoU on the 14 known classes drops by XX.X mIoU {[}TBD-RQ2{]}
versus the 19-class fully-supervised baseline. A robustness 16-known /
3-unknown split (drop only \texttt{bicyclist}, \texttt{motorcyclist},
\texttt{other-vehicle}) reports AUROC = XX.X {[}TBD-RQ2{]}; a reverse
cross-domain check trained on the SemanticKITTI 14-known split and
evaluated on nuScenes-LiDARSeg classes that have no SemanticKITTI
analogue (\texttt{construction\_vehicle}, \texttt{barrier},
\texttt{traffic\_cone}, \texttt{pushable\_pullable}) reports AUROC =
XX.X {[}TBD-RQ2{]}. Figure 5 {[}TBD-after-exp{]} shows the bimodal
vacuity histogram on a representative val frame: known-class voxels
cluster at \(u_v \in (0, 0.2]\), withheld-class voxels at
\(u_v \in [0.6, 1.0]\), with the cross-over bin carrying \(< 5\%\) of
voxel mass --- the visual demonstration of H2.

\textbf{TABLE V. RQ2 --- OPEN-SET VACUITY AS OOD SCORE; THREE SPLITS ×
THREE METRICS × FOUR METHODS.}

{\def\LTcaptype{none} % do not increment counter
\begin{longtable}[]{@{}lllllll@{}}
\toprule\noalign{}
Method & 14/5 AUROC & 14/5 AUPR & 14/5 mIoU on Known & 16/3 AUROC & 16/3
AUPR & nuScenes Reverse AUROC \\
\midrule\noalign{}
\endhead
\bottomrule\noalign{}
\endlastfoot
R2-reimpl {[}1{]} & XX.X & XX.X & XX.X & XX.X & XX.X & XX.X \\
ConvBKI {[}4{]} & XX.X & XX.X & XX.X & XX.X & XX.X & XX.X \\
Clio-LiDAR-stub {[}18{]} & XX.X & XX.X & XX.X & XX.X & XX.X & XX.X \\
\textbf{EvidLife-Map} & \textbf{XX.X} & \textbf{XX.X} & \textbf{XX.X} &
\textbf{XX.X} & \textbf{XX.X} & \textbf{XX.X} \\
\end{longtable}
}

{[}TBD-after-exp{]} One paragraph contrasting our per-voxel Dirichlet
vacuity against ConvBKI\textquotesingle s kernel-smoothed variance: we
expect ConvBKI\textquotesingle s smoothing to \emph{suppress} per-voxel
epistemic spikes on isolated unknown points, lowering its OOD AUROC even
when its closed-set mIoU is competitive --- the entire reason §III.B
drops the spatial kernel.

\subsection{IV.E Downstream Traversability --- RQ3 (SemanticSpray
Closed-Loop
Replay)}\label{ive-downstream-traversability--rq3-semanticspray-closed-loop-replay}

On the SemanticSpray {[}21{]} passive-replay traversability evaluation,
propagating per-voxel vacuity into a downstream traversability head
yields XX.X percentage-point absolute increase in safe-region recall and
a XX\% relative reduction in false-traversable rate {[}TBD-RQ3{]} versus
an R2-style hard-rule baseline at matched coverage. The deferral rate
(fraction of voxels flagged "do not commit to traversable") increases by
XX pp {[}TBD-RQ3{]} in compensation, an expected and desired behaviour
for an open-set-aware traversability layer that prefers caution to false
confidence on visually-ambiguous wet-road pixels.

\textbf{TABLE VI. RQ3 --- TRAVERSABILITY ON SEMANTICSPRAY {[}21{]}
PASSIVE REPLAY.}

{\def\LTcaptype{none} % do not increment counter
\begin{longtable}[]{@{}lllll@{}}
\toprule\noalign{}
Method & Safe-Region Recall (\%) & False-Traversable Rate (\%) &
Deferral Rate (\%) & Coverage (\%) \\
\midrule\noalign{}
\endhead
\bottomrule\noalign{}
\endlastfoot
nvblox-vanilla + hard rule {[}3{]} & XX.X & XX.X & XX.X & XX.X \\
R2-style hard rule {[}1{]} & XX.X & XX.X & XX.X & XX.X \\
\textbf{EvidLife-Map (vacuity-gated)} & \textbf{XX.X} & \textbf{XX.X} &
\textbf{XX.X} & \textbf{XX.X} \\
\end{longtable}
}

\subsection{IV.F Lifelong Multi-Session --- RQ4 (KITTI-360 Revisit
Pairs)}\label{ivf-lifelong-multi-session--rq4-kitti-360-revisit-pairs}

Across XX revisit pairs {[}TBD-RQ4{]} constructed by odometry-overlap on
KITTI-360 sequences 00, 02, 04, 05, 06, 07, 09, 10, EvidLife-Map
achieves \textbf{stale-voxel removal precision} XX\% at recall XX\%
{[}TBD-RQ4{]} over the second-pass trajectory after the M3 decay (Eq.
12) has acted, \textbf{ECE drift across sessions} of XX× {[}TBD-RQ4{]}
the single-session ECE (target \(\le 1.5\times\)), and a
\textbf{multi-session memory footprint} of XX× {[}TBD-RQ4{]} the
single-session footprint over the five revisited pairs (target
\(\le 1.7\times\), sub-linear in session count thanks to the voxel-hash
deduplication described in §III.F). Figure 6 {[}TBD-after-exp{]} shows
per-pair stale-voxel PR curves; the per-pair area-under-curve degrades
gracefully with revisit gap length, confirming that the vacuity-driven
decay schedule does not catastrophically forget under long absences.

\textbf{TABLE VII. RQ4 --- LIFELONG MULTI-SESSION ON KITTI-360 REVISIT
PAIRS.}

{\def\LTcaptype{none} % do not increment counter
\begin{longtable}[]{@{}lllllll@{}}
\toprule\noalign{}
Revisit Pair & Stale-Voxel P (\%) & Stale-Voxel R (\%) & F1 (\%) &
Single-Session ECE & Multi-Session ECE & Memory Growth (\(\times\)) \\
\midrule\noalign{}
\endhead
\bottomrule\noalign{}
\endlastfoot
seq 00 --- pair 1 & XX & XX & XX & XX.X & XX.X & XX.X \\
seq 02 --- pair 2 & XX & XX & XX & XX.X & XX.X & XX.X \\
seq 04 --- pair 3 & XX & XX & XX & XX.X & XX.X & XX.X \\
seq 05 --- pair 4 & XX & XX & XX & XX.X & XX.X & XX.X \\
seq 06 --- pair 5 & XX & XX & XX & XX.X & XX.X & XX.X \\
\textbf{Mean} & \textbf{XX} & \textbf{XX} & \textbf{XX} & \textbf{XX.X}
& \textbf{XX.X} & \textbf{XX.X} \\
\end{longtable}
}

\subsection{IV.G Dynamic Scene Comparison vs Khronos --- RQ5
(SemanticKITTI Dynamic
Split)}\label{ivg-dynamic-scene-comparison-vs-khronos--rq5-semantickitti-dynamic-split}

On the SemanticKITTI dynamic-frame subset (frames from seq 00, 04, 05,
07 with \(\ge 5\%\) \texttt{moving-*} GT points), EvidLife-Map achieves
dynamic-object mIoU XX.X (Khronos {[}6{]} published / reproduced XX.X),
static-region recall XX.X pp (\$\textbackslash Delta = \$ XX pp vs
Khronos), and end-to-end latency XX.X ms on Jetson Orin NX
{[}TBD-RQ5{]}. The framing established in §III.D (implicit
vacuity-instability vs explicit 4D factoriser) is the basis of the
comparison: we do not claim that the implicit channel dominates the
explicit one on dynamic-object mIoU; we claim it is competitive enough
to make Khronos\textquotesingle s full short-/long-term machinery a
deployment-cost trade-off rather than a capability necessity in our
operating regime. Figure 5 {[}TBD-after-exp{]} (the dynamic version,
distinct from the §IV.D vacuity histogram) shows a side-by-side
qualitative comparison on a representative dynamic frame.

\textbf{TABLE VIII. RQ5 --- DYNAMIC-SCENE HEAD-TO-HEAD ON THE
SEMANTICKITTI DYNAMIC SPLIT.}

{\def\LTcaptype{none} % do not increment counter
\begin{longtable}[]{@{}llllll@{}}
\toprule\noalign{}
Method & Dynamic-Object mIoU & Static-Region Recall (\%) & E2E Latency
(ms, 4060) & E2E Latency (ms, Orin NX) & GPU Memory (GB) \\
\midrule\noalign{}
\endhead
\bottomrule\noalign{}
\endlastfoot
R2-reimpl {[}1{]} (lower bound) & XX.X & XX.X & XX.X & XX.X & XX.X \\
Khronos {[}6{]} (reproduced / publ.) & XX.X & XX.X & XX.X & not designed
for Orin NX & XX.X \\
\textbf{EvidLife-Map} & \textbf{XX.X} & \textbf{XX.X} & \textbf{XX.X} &
\textbf{XX.X} & \textbf{XX.X} \\
\end{longtable}
}

\subsection{IV.H Ablations}\label{ivh-ablations}

We ablate the six firm pipeline choices of research\_plan v3 §4.4 (A-1
through A-5 plus A-7) and one stretch choice (A-6, voxel size) reported
only if W14 gates are green. Each row of Tab IX disables exactly one
design choice; the column triple (mIoU / ECE / open-set AUROC) isolates
the effect on the three reviewer-credible RQs simultaneously.

\textbf{TABLE IX. ABLATIONS (1 × RTX 4060, 16 GB).}

{\def\LTcaptype{none} % do not increment counter
\begin{longtable}[]{@{}lllllll@{}}
\toprule\noalign{}
\# & Ablation & SK mIoU & SK ECE & Open-Set AUROC & KITTI-360 Stale P
(\%) & Orin NX Hz \\
\midrule\noalign{}
\endhead
\bottomrule\noalign{}
\endlastfoot
Full & EvidLife-Map (full) & XX.X & XX.X & XX.X & XX & XX.X \\
A-1 & \(-\) Dirichlet head (use argmax-Bayes) & XX.X & XX.X & XX.X & XX
& XX.X \\
A-2 & \(-\) vacuity-coupled \(\tau\) (fixed-\(\tau\) decay) & XX.X &
XX.X & XX.X & XX & XX.X \\
A-3 & \(-\) entropy channel in descriptor (use \(h_{\text{class}}\)
only) & XX.X & XX.X & XX.X & XX & XX.X \\
A-4 & \(-\) vacuity-aware traversability (hard threshold) & XX.X & XX.X
& XX.X & XX & XX.X \\
A-5 & \(-\) conjugate Dirichlet inter-session fusion (naive overwrite) &
XX.X & XX.X & XX.X & XX & XX.X \\
A-7 & \(-\) dedicated unknown channel (vacuity only, \(C\)=14) & XX.X &
XX.X & XX.X & XX & XX.X \\
A-6 (stretch) & voxel size \(0.10\,\text{m}\) vs \(0.25\,\text{m}\) &
XX.X & XX.X & XX.X & XX & XX.X \\
\end{longtable}
}

{[}TBD-after-exp{]} One paragraph naming which ablation is
"load-bearing" (likely A-2 for stale-voxel precision, A-3 for
loop-closure precision, A-7 for open-set AUROC) and which is small
enough to demote to supplementary in a camera-ready trim if §IV.G
overflows.

\subsection{IV.I Runtime, Memory, and Jetson Deployment + Demo
Video}\label{ivi-runtime-memory-and-jetson-deployment--demo-video}

On a single RTX 4060 (16 GB) the end-to-end pipeline runs at XX.X Hz
{[}TBD-RQ-deploy{]} over a SemanticKITTI seq 08 replay; on the Jetson
Orin NX (16 GB) with a TensorRT FP16 build of the M1 head and a CPU
implementation of M2 loop closure, the pipeline runs at XX.X Hz
{[}TBD-RQ-deploy{]} over the same replay, peak VRAM XX.X GB. A 30-second
supplementary video \texttt{evidlife\_demo.mp4} shows EvidLife-Map
running live on the Orin NX with the metric-semantic voxel map, the
traversability overlay (RQ3 output), and a vacuity heat-map overlay
visible in real time at native 10 Hz over 300 frames; capture and
overlay scripts ship with the code release.

\textbf{TABLE X. RUNTIME, MEMORY, AND JETSON DEPLOYMENT.}

{\def\LTcaptype{none} % do not increment counter
\begin{longtable}[]{@{}lll@{}}
\toprule\noalign{}
Metric & RTX 4060 (16 GB) & Jetson Orin NX (16 GB) \\
\midrule\noalign{}
\endhead
\bottomrule\noalign{}
\endlastfoot
Mean end-to-end latency (ms) & XX.X & XX.X \\
99-pct latency (ms) & XX.X & XX.X \\
Throughput (Hz) & XX.X & XX.X \\
Peak GPU memory (GB) & XX.X & XX.X \\
Per-voxel memory (B) at \(C=19\) & 180 & 180 \\
Per-block memory (kB) at \(8^3\) & 44 & 44 \\
\end{longtable}
}

\begin{center}\rule{0.5\linewidth}{0.5pt}\end{center}

\section{V. Discussion}\label{v-discussion}

\subsection{V.A Limitations}\label{va-limitations}

Five limitations are stated honestly to keep the §I.B thesis falsifiable
at every revision. The first three are pre-registered scope decisions
from research\_plan v3, the last two are surfaced by the §IV.0
preliminary verification campaign.

\textbf{(i) RGB head omitted.} The system is LiDAR-only at deployment;
we use RGB only as a training-time prior in the Cylinder3D pretrain, and
we make no claim about photo-realistic appearance, weather robustness,
or open-vocabulary text retrieval. The Gaussian-splatting strand
exemplified by GS-LIVO {[}20{]} and the open-vocabulary scene-graph
strand exemplified by HOV-SG {[}19{]} are therefore methodologically
orthogonal future work, not in-scope contenders here.

\textbf{(ii) Static voxel grid, no Gaussian / NeRF representation.} We
keep the nvblox {[}3{]} voxel grid because the M3 decay (Eq. 12)
requires a per-voxel state with bounded memory and well-defined cell
boundaries for the conjugate update; a continuous-field representation
would need a different decay primitive that we have not derived.

\textbf{(iii) Per-voxel \(\alpha\) memory scaling at \(C \ge 100\).} The
\((C+1) \times 4\,\text{B}\) growth of the per-voxel record becomes the
dominant memory cost when the closed-set taxonomy is much larger than
the 19-class SemanticKITTI setting; the open-vocabulary regime where
\(C\) is the size of a CLIP token vocabulary therefore needs either a
low-rank approximation of \(\alpha_v\) or a per-instance Bernoulli-like
collapse along OpenVox {[}5{]} lines. We deliberately do not claim
either extension in this paper.

\textbf{(iv) Backbone path declaration --- Path 4 (PointNet2Lite\_v2) is
the pre-registered §IV.0 backbone; Cylinder3D extension is journal-track
future work.} The planned Cylinder3D backbone is gated on a spconv 1.x
source build (supplementary \texttt{W2-1\_status.md}: 0 missing keys
after weight permutation + 36 \texttt{.features} →
\texttt{.replace\_feature} patches + 22 kernel-shape-disambiguated
\texttt{indice\_key} patches, but spconv 1.x → 2.x kernel-index ordering
value-semantic gap produces NaN logits). Of the three §3.9 unblock paths
originally considered, \textbf{we pre-register Path 4 (extended
PointNet2Lite\_v2, 0.49 M params)} as this paper\textquotesingle s
backbone, with the spconv 1.x source build (Path 1) and Cylinder3D-class
re-implementation deliberately deferred to a journal extension. The
pre-registration is grounded in three findings from the W3 campaign: (a)
the §IV.0 W3-R single-checkpoint joint-serving result (mIoU 30.31 \% on
14-known + vacuity AUROC 0.706 + M3 stale-voxel F1 0.76) is a
deployment-grade demonstration of the §I.B thesis without Cylinder3D;
(b) PVKD was inspected and found to inherit the same spconv 1.x
dependency (\texttt{backbone\_path\_findings.md}), so Path 2 is ruled
out; (c) Path 1 source-build attempts on torch 2.11 + cu130 have a 60 \%
expected failure rate (R-21 in research\_plan v4), so deferring it
preserves the §IV submission window. Reviewers should evaluate the §IV.0
contribution at the PointNet2Lite\_v2 tier as the \emph{intended}
preliminary result rather than a stand-in for a missing Cylinder3D; the
Cylinder3D-extrapolation language previously in §VI has been removed.

\textbf{(v) Preliminary mIoU absolute values are neighbourhood-bound at
the PointNet capacity tier.} Multi-sequence audit (§IV.0 revised)
measures M1+kNN at 17.92 \% mIoU on the official SemKITTI split, still
far below the ConvBKI {[}4{]} published 77.7 \%. Two factors combine:
(a) the PointNet capacity tier (0.21--0.28 M params) sits ≈ 200× below
Cylinder3D (55 M), and (b) even the lite variant uses only a single
kNN-pool layer, while Cylinder3D\textquotesingle s sparse-conv
cylindrical decoder propagates context across the entire voxel volume.
The §IV.0 multi-seq result \emph{isolates} the per-method gain: at
matched backbone tier, M1 (EDL) consistently improves ECE by 4--7× and,
with the kNN extension, also improves mIoU by 7.8× vs R2 (CE) --- the
genuine method-level signal independent of absolute mIoU. The full-scale
§IV.C numbers replace these preliminary cells once limitation (iv) is
unblocked.

\subsection{V.B Failure Modes}\label{vb-failure-modes}

Four failure modes are documented; (a)-(c) are pre-registered from the
planning audit and (d) is observed in the §IV.0 preliminary RQ2 run.

\textbf{(a) Confidently-wrong evidence accumulation.} Persistent
mirror-surface returns or other systematic LiDAR ghosting cause evidence
to accumulate with low entropy in a wrong class, keeping vacuity
artificially low and preventing M3 from triggering decay. This is a
known evidential-deep-learning failure mode inherited from Sensoy 2018,
and it also produces low-vacuity false negatives in the RQ2 open-set
evaluation for unknown classes that visually resemble a known class
(e.g., a truck withheld from training that gets confident \texttt{car}
evidence on every frame).

\textbf{(b) Non-overlapping class distributions across sessions.} When
two sessions of the same area capture non-overlapping class sets (e.g.,
one before construction, one after), the M2 class-histogram channel of
the descriptor collapses to near-zero cosine similarity and the matching
degenerates to the entropy channel only; loop-closure precision degrades
by approximately XX percentage points {[}TBD-after-exp{]} in such cases.

\textbf{(c) Taxonomically-close withheld classes.} On the 16/3
robustness open-set split (withhold only \texttt{bicyclist},
\texttt{motorcyclist}, \texttt{other-vehicle}) the AUROC degrades by
XX.X {[}TBD-after-exp{]} relative to the 14/5 primary split, because the
withheld classes have close training analogues (\texttt{bicycle},
\texttt{motorcycle}, \texttt{car}) whose evidence channels capture most
of the unknown points\textquotesingle{} mass.

\textbf{(d) EDL post-epoch-1 vacuity collapse --- historical trade-off
resolved by W3-R single-schedule open-set training.} Earlier W3
ablations exposed a dimension-specific KL trade-off: at the
closed-set-trained W3-H/J/M tier, KL warm-restart maximised mIoU (W3-M
23.32 \% at \(\lambda_t\) restart) but strictly hurt vacuity AUROC (W3-Q
0.776 \textless{} W3-P 0.781 linear-KL; cycle-restart valleys reached
AUROC 0.43--0.58 because MSE-only loss collapsed vacuity for all points
including unknowns the head had never seen). This led the v2 draft to
maintain \emph{two} checkpoints with \emph{two} KL schedules and select
per target metric --- a model-selection oracle no real deployment can
access (round-1 R1 W2, R3 W2, R4 CRITICAL \#1). \textbf{The W3-R
single-checkpoint joint-serving experiment (§IV.0; supplementary
\texttt{w3r\_joint\_serving.json}) resolves the trade-off}: under a
\emph{single} linear-KL schedule with open-set training (14-known +
5-unknown masked to ignore-index), one selected checkpoint (ep 12, best
mIoU subject to AUROC ≥ 0.70 floor) achieves mIoU 30.31 \% + AUROC 0.706
+ ECE 0.320 + M3 vacuity F1 0.76 simultaneously, demonstrating that the
integration thesis IS deployable from a single shipped network. The
original warm-restart-vs-linear trade-off is now reported as a
\emph{closed-set-training-protocol artefact} that the open-set training
protocol resolves; the dimension-specific failure mode this footnote
previously named is empirically retired by W3-R.

Figure 6 {[}TBD-after-exp{]} gives one qualitative panel per failure
mode.

\subsection{V.C Honest Negative
Findings}\label{vc-honest-negative-findings}

\textbf{The pre-registered RQ2 G-5 gate is approached but not cleared at
multi-sequence scale; W3-C\textquotesingle s 0.808 was spatially
inflated.} The W3-C/N/O/P audit chain (§IV.D table) confirms (a) the
W3-C 0.808 came from single-seq spatial-adjacency inflation, not from a
genuine method advantage at single-seq scale, and (b) the honest
multi-seq operating point is \textbf{0.781} (W3-P, lite\_v2 backbone,
16/3 robustness split). The §I.B "one vacuity, three jobs" thesis
therefore stays intact in \textbf{direction} --- vacuity discriminates
unknowns above chance with a clear positive margin (+0.28 above random)
--- but the \emph{strength} of the leg (i) claim is reframed: vacuity is
a useful OOD score, not a maximum-quality one. The pre-registered C1
fallback ("two jobs verified, leg (i) protocol-dependent") is therefore
the honest §I.B framing for the multi-seq operating regime; the W3-C
0.808 stays in the paper only as a documented operating point with its
protocol caveat. The "Cylinder3D regime" hypothesis remains testable: at
the full backbone scale (Limitation V.A iv) we expect AUROC to lift
toward 0.85+ because more spatial frequency information is available to
discriminate; W3-N → W3-P shows the +0.04 gain a deeper backbone already
gives at our scale.

\textbf{Preliminary absolute mIoU does not yet match published
ConvBKI/S-BKI numbers.} Section IV.0 reports an M1+kNN mIoU of 17.92 \%
on the official multi-sequence split versus published targets of 77.7 \%
(ConvBKI {[}4{]}, KITTI seq 15) and 51.3 \% (S-BKI {[}2{]},
SemanticKITTI test). This \textasciitilde{} 50--60 pp absolute gap is
dominated by the joint capacity-and-context gap to Cylinder3D
(Limitation iv-v); the comparative ordering at matched backbone (M1
\textgreater{} R2 by 4--7× on ECE, and by 7.8× on mIoU once kNN context
is enabled) is the genuine method-level signal we currently claim.
Full-scale absolute numbers replace the §IV.0 placeholders once
Limitation iv is unblocked.

\textbf{Khronos head-to-head pre-registered fallback.} If the RQ5
head-to-head against Khronos {[}6{]} in §IV.G shows a dynamic-object
mIoU gap of more than 3 against Khronos\textquotesingle s published
numbers and the latency advantage on Orin NX is less than \(2\times\),
we will explicitly acknowledge it and reframe the win as a
latency-integration trade-off rather than a methodological replacement
--- closing the audit E6 ("over-claiming") risk before a reviewer raises
it.

\begin{center}\rule{0.5\linewidth}{0.5pt}\end{center}

\section{VI. Conclusion}\label{vi-conclusion}

We have presented EvidLife-Map, a LiDAR-only online metric-semantic
mapping system whose closed-form Dirichlet-evidential per-voxel
posterior exposes two complementary scalars --- vacuity and dissonance
--- under one conjugate update primitive, allocated per-job across three
downstream consumers: vacuity drives open-set / OOD detection (M1) and
lifelong voxel decay (M3); dissonance drives the loop-closure descriptor
entropy channel (M2). The three modules share one per-voxel state
\(\alpha_v\) and one arithmetic primitive (conjugate Dirichlet addition
/ multiplicative pull-toward-prior). At the pre-registered Path-4
backbone (§V.A iv: PointNet2Lite\_v2, 0.49 M params), the W3-R
single-checkpoint joint-serving experiment delivers all four headline
metrics from one shipped network: \textbf{mIoU 30.31 \% on 14 known
classes, vacuity AUROC 0.706 on the 14-known / 5-unknown open-set split,
ECE 0.320, and M3 stale-voxel F1 0.76 at threshold 0.5} --- the first
reported demonstration that one shipped EDL-MSM model serves closed-set
semantic mapping, open-set OOD discrimination, lifelong voxel decay, and
loop-closure descriptor simultaneously. Complementary numbers at the
closed-set-trained tier (W3-M 23.32 \% mIoU / 0.125 ECE) and the
matched-protocol R2 baseline (W3-S CE-lite\_v2 22.54 \% / 0.089 ECE)
jointly close the round-1 R1 W1 protocol-symmetry critique. The
Cylinder3D-class extension (§V.A iv) is explicitly deferred to
journal-track future work; the §IV.0 contribution is the \emph{intended}
preliminary result at PointNet2Lite\_v2 tier, not a stand-in for a
missing larger backbone. Full code, multi-sequence training scripts,
per-class IoU dumps, the W3-A through W3-UVW / R / S audit trail, and
the simulated round-1 peer-review reports are released on the project
repository. Future work explicitly out of scope for this paper but
enabled by the same per-voxel evidential substrate includes (i) a
SAM-based open-vocabulary head feeding evidence into the same Dirichlet
update, (ii) a 3D Gaussian-splatting variant of the same vacuity-decay
coupling on a non-voxel substrate (the GS-LIVO {[}20{]} direction), and
(iii) a 4D-radar fusion branch that supplies an evidence channel
orthogonal to LiDAR under adverse weather --- the latter restoring the
v1 multi-modal ambition that the v3 deployment-on-4060 scope
deliberately excluded.

\begin{center}\rule{0.5\linewidth}{0.5pt}\end{center}

\section{References}\label{references}

\emph{Numeric → bibkey map, indexed by first appearance in §I--§VI. All
bibkeys reference entries of
\texttt{D:\textbackslash{}\_7\_sci\textbackslash{}semantic\_mapping\textbackslash{}\_new\_paper\textbackslash{}refs\textbackslash{}refs.bib}
(60 entries; this paper cites 30/60).}

{[}1{]} J. Jiao, R. Geng, Y. Li, R. Xin, B. Yang, J. Wu, L. Wang, M.
Liu, R. Fan, and D. Kanoulas, "Real-Time Metric-Semantic Mapping for
Autonomous Navigation in Outdoor Environments," \emph{IEEE Trans. Autom.
Sci. Eng.}, 2024. \emph{{[}bibkey: jiao2024r2{]}}

{[}2{]} L. Gan, R. Zhang, J. W. Grizzle, R. M. Eustice, and M. Ghaffari,
"Bayesian Spatial Kernel Smoothing for Scalable Dense Semantic Mapping,"
\emph{IEEE Robot. Autom. Lett.}, 2020. \emph{{[}bibkey: gan2020sbki{]}}

{[}3{]} A. Millane, H. Oleynikova, E. Wirbel, R. Steiner, V. Ramasamy,
D. Tingdahl, and R. Siegwart, "nvblox: GPU-Accelerated Incremental
Signed Distance Field Mapping," in \emph{Proc. IEEE Int. Conf. Robot.
Autom. (ICRA)}, 2024. \emph{{[}bibkey: millane2024nvblox{]}}

{[}4{]} J. Wilson, J. Song, Y. Fu, A. Zhang, A. Capodieci, P. Jayakumar,
K. Barton, and M. Ghaffari, "ConvBKI: Real-Time Probabilistic Semantic
Mapping Network with Quantifiable Uncertainty," \emph{IEEE Trans.
Robot.}, 2024. \emph{{[}bibkey: wilson2024convbki{]}}

{[}5{]} Y. Deng, B. Yao, Y. Tang, Y. Yang, and Y. Yue, "OpenVox:
Real-time Instance-level Open-vocabulary Probabilistic Voxel
Representation," \emph{arXiv:2502.16528}, 2025. \emph{{[}bibkey:
deng2025openvox{]}}

{[}6{]} L. Schmid, M. Abate, Y. Chang, and L. Carlone, "Khronos: A
Unified Approach for Spatio-Temporal Metric-Semantic SLAM in Dynamic
Environments," in \emph{Proc. Robot.: Sci. Syst. (RSS)}, 2024.
\emph{{[}bibkey: schmid2024khronos{]}}

{[}7{]} J. Wilson, R. Xu, Y. Sun, P. Ewen, M. Zhu, K. Barton, and M.
Ghaffari, "LatentBKI: Open-Dictionary Continuous Mapping in
Visual-Language Latent Spaces with Quantifiable Uncertainty,"
\emph{arXiv:2410.11783}, 2024. \emph{{[}bibkey: wilson2024latentbki{]}}

{[}8{]} L. Schmid, J. Delmerico, J. L. Chung, M. Magnusson, J. Nieto,
and R. Siegwart, "Panoptic Multi-TSDFs: a Flexible Representation for
Online Multi-resolution Volumetric Mapping and Long-term Dynamic Scene
Consistency," in \emph{Proc. IEEE Int. Conf. Robot. Autom. (ICRA)},
2022. \emph{{[}bibkey: schmid2022panopticmultitsdfs{]}}

{[}9{]} A. Rosinol, A. Violette, M. Abate, N. Hughes, Y. Chang, J. Shi,
A. Gupta, and L. Carlone, "Kimera: From SLAM to Spatial Perception with
3D Dynamic Scene Graphs," \emph{Int. J. Robot. Res.}, 2021.
\emph{{[}bibkey: rosinol2021kimera{]}}

{[}10{]} N. Hughes, Y. Chang, and L. Carlone, "Hydra: A Real-time
Spatial Perception System for 3D Scene Graph Construction and
Optimization," in \emph{Proc. Robot.: Sci. Syst. (RSS)}, 2022.
\emph{{[}bibkey: hughes2022hydra{]}}

{[}11{]} H. Oleynikova, Z. Taylor, M. Fehr, R. Siegwart, and J. Nieto,
"Voxblox: Incremental 3D Euclidean Signed Distance Fields for On-Board
MAV Planning," in \emph{Proc. IEEE/RSJ Int. Conf. Intell. Robots Syst.
(IROS)}, 2017. \emph{{[}bibkey: oleynikova2017voxblox{]}}

{[}12{]} A. Hornung, K. M. Wurm, M. Bennewitz, C. Stachniss, and W.
Burgard, "OctoMap: An Efficient Probabilistic 3D Mapping Framework Based
on Octrees," \emph{Auton. Robots}, 2013. \emph{{[}bibkey:
hornung2013octomap{]}}

{[}13{]} M. Grinvald, F. Furrer, T. Novkovic, J. J. Chung, C. Cadena, R.
Siegwart, and J. Nieto, "Volumetric Instance-Aware Semantic Mapping and
3D Object Discovery," \emph{IEEE Robot. Autom. Lett.}, 2019.
\emph{{[}bibkey: grinvald2019voxbloxpp{]}}

{[}14{]} G. Narita, T. Seno, T. Ishikawa, and Y. Kaji, "PanopticFusion:
Online Volumetric Semantic Mapping at the Level of Stuff and Things," in
\emph{Proc. IEEE/RSJ Int. Conf. Intell. Robots Syst. (IROS)}, 2019.
\emph{{[}bibkey: narita2019panopticfusion{]}}

{[}15{]} Y. Pan, Y. Kompis, L. Bartolomei, R. Mascaro, C. Stachniss, and
M. Chli, "Voxfield: Non-Projective Signed Distance Fields for Online
Planning and 3D Reconstruction," in \emph{Proc. IEEE/RSJ Int. Conf.
Intell. Robots Syst. (IROS)}, 2022. \emph{{[}bibkey: pan2022voxfield{]}}

{[}16{]} Y. Tian, Y. Chang, F. H. Arias, C. Nieto-Granda, J. P. How, and
L. Carlone, "Kimera-Multi: Robust, Distributed, Dense Metric-Semantic
SLAM for Multi-Robot Systems," \emph{IEEE Trans. Robot.}, 2022.
\emph{{[}bibkey: tian2022kimeramulti{]}}

{[}17{]} Y. Chang, N. Hughes, A. Ray, and L. Carlone, "Hydra-Multi:
Collaborative Online Construction of 3D Scene Graphs with Multi-Robot
Teams," in \emph{Proc. IEEE/RSJ Int. Conf. Intell. Robots Syst. (IROS)},
2023. \emph{{[}bibkey: chang2023hydramulti{]}}

{[}18{]} D. Maggio, Y. Chang, N. Hughes, M. Trang, D. Griffith, C.
Dougherty, E. Cristofalo, L. Schmid, and L. Carlone, "Clio: Real-Time
Task-Driven Open-Set 3D Scene Graphs," \emph{IEEE Robot. Autom. Lett.},
2024. \emph{{[}bibkey: maggio2024clio{]}}

{[}19{]} A. Werby, C. Huang, M. Büchner, A. Valada, and W. Burgard,
"HOV-SG: Hierarchical Open-Vocabulary 3D Scene Graphs for
Language-Grounded Robot Navigation," in \emph{Proc. Robot.: Sci. Syst.
(RSS)}, 2024. \emph{{[}bibkey: werby2024hovsg{]}}

{[}20{]} Z. Hong, X. Zheng, W. Ding, S. Shen, and Y. Zhou, "GS-LIVO:
Real-Time LiDAR, Inertial, and Visual Multisensor Fused Odometry with
Gaussian Mapping," \emph{IEEE Trans. Robot.}, 2025. \emph{{[}bibkey:
hong2025gslivo{]}}

{[}21{]} A. Piroli, V. Dallabetta, J. Kopp, M. Walessa, D. Meissner, and
K. Dietmayer, "Energy-Based Detection of Adverse Weather Effects in
LiDAR Data," \emph{IEEE Robot. Autom. Lett.}, 2023. \emph{{[}bibkey:
piroli\_2023\_semanticspray{]}}

{[}22{]} K. Yamazaki, T. Hanyu, K. Vo, T. Pham, M. Tran, G. Doretto, A.
Nguyen, and N. Le, "Open-Fusion: Real-time Open-Vocabulary 3D Mapping
and Queryable Scene Representation," in \emph{Proc. IEEE Int. Conf.
Robot. Autom. (ICRA)}, 2024. \emph{{[}bibkey: yamazaki2024openfusion{]}}

{[}23{]} A. Sheppard, P. Ewen, J. Wilson, A. V. Sethuraman, B. Adewole,
A. Li, Y. Chen, R. Vasudevan, and K. A. Skinner, "SLIM-VDB: A Real-Time
3D Probabilistic Semantic Mapping Framework," \emph{IEEE Robot. Autom.
Lett.}, 2025. \emph{{[}bibkey: sheppard2025slimvdb{]}}

{[}24{]} J. Fu, C. Lin, Y. Taguchi, A. Cohen, Y. Zhang, S. Mylabathula,
and J. J. Leonard, "PlaneSDF-Based Change Detection for Long-term Dense
Mapping," \emph{IEEE Robot. Autom. Lett.}, 2022. \emph{{[}bibkey:
fu\_2022\_planesdf{]}}

{[}25{]} M. A. Vega-Torres, A. Braun, and A. Borrmann, "SLAM2REF:
Advancing Long-Term Mapping with 3D LiDAR and Reference Map Integration
for Precise 6-DoF Trajectory Estimation and Map Extension,"
\emph{Constr. Robot.}, 2024. \emph{{[}bibkey:
vegatorres\_2024\_slam2ref{]}}

{[}26{]} L. Li, X. Kong, X. Zhao, W. Li, F. Wen, H. Zhang, and Y. Liu,
"SA-LOAM: Semantic-aided LiDAR SLAM with Loop Closure," in \emph{Proc.
IEEE Int. Conf. Robot. Autom. (ICRA)}, 2021. \emph{{[}bibkey:
li\_2021\_saloam{]}}

{[}27{]} L. Yang, R. Mascaro, I. Alzugaray, S. M. Prakhya, M. Karrer, Z.
Liu, and M. Chli, "LiDAR Loop Closure Detection using Semantic Graphs
with Graph Attention Networks," \emph{J. Intell. Robot. Syst.}, 2025.
\emph{{[}bibkey: yang\_2025\_semanticloop{]}}

{[}28{]} J. L. Matez-Bandera, D. Fernandez-Chaves, J.-R. Ruiz-Sarmiento,
J. Monroy, N. Petkov, and J. Gonzalez-Jimenez, "LTC-Mapping: Enhancing
Long-Term Consistency of Object-Oriented Semantic Maps in Robotics,"
\emph{Sensors}, 2022. \emph{{[}bibkey: matezbandera2022ltcmapping{]}}

{[}29{]} Y. Miao, I. Armeni, M. Pollefeys, and D. Barath, "Volumetric
Semantically Consistent 3D Panoptic Mapping," \emph{arXiv:2309.14737},
2023. \emph{{[}bibkey: miao2023volumetric{]}}

{[}30{]} S. Zhu, R. Qin, G. Wang, J. Liu, and H. Wang, "SemGauss-SLAM:
Dense Semantic Gaussian Splatting SLAM," in \emph{Proc. IEEE/RSJ Int.
Conf. Intell. Robots Syst. (IROS)}, 2025. \emph{{[}bibkey:
zhu2025semgaussslam{]}}

\begin{center}\rule{0.5\linewidth}{0.5pt}\end{center}

\section{Appendix}\label{appendix}

\textbf{Acknowledgments.} {[}TBD{]}

\textbf{Author Contributions.} {[}TBD{]}

\end{document}
